% =========================================================================
% Width-Bounded Heuristics for Symbolic Search
% Draft for ICAPS 2027. Compile with pdflatex; port to icaps27.sty later.
% Experiments section reports results from the SymK implementation.
% =========================================================================
\documentclass[letterpaper]{article}
\usepackage[margin=1in]{geometry}
\usepackage{amsmath,amssymb,amsthm}
\usepackage{natbib}
\usepackage{booktabs}
\usepackage{microtype}

% ------------------------------- macros ---------------------------------
\newcommand{\inlinecite}[1]{\citet{#1}}
\newcommand{\task}{\ensuremath{\Pi}}
\newcommand{\vars}{\ensuremath{\mathcal{V}}}
\newcommand{\ops}{\ensuremath{\mathcal{O}}}
\newcommand{\states}{\ensuremath{\mathcal{S}}}
\newcommand{\sinit}{\ensuremath{s_{\mathrm{I}}}}
\newcommand{\sgoal}{\ensuremath{s_\star}}
\newcommand{\pre}{\ensuremath{\mathit{pre}}}
\newcommand{\eff}{\ensuremath{\mathit{eff}}}
\newcommand{\cost}{\ensuremath{\mathit{cost}}}
\newcommand{\hstar}{\ensuremath{h^{*}}}
\newcommand{\hblind}{\ensuremath{h^{0}}}
\newcommand{\hpot}{\ensuremath{h^{\text{pot}}}}
\newcommand{\gstar}{\ensuremath{g^{*}}}
\newcommand{\copt}{\ensuremath{C^{*}}}
\newcommand{\bdd}[1]{\ensuremath{B_{#1}}}
\newcommand{\bddsize}[1]{\ensuremath{|B_{#1}|}}
\newcommand{\cofcount}{\ensuremath{Q}}
\newcommand{\width}{\ensuremath{\mathsf{w}}}
\newcommand{\values}{\ensuremath{V}}
\newcommand{\level}[2]{\ensuremath{H_{#2}}}
\newcommand{\effort}{\ensuremath{\mathit{effort}}}
\newcommand{\sub}[2]{\ensuremath{\mathrm{cof}_{#1}(#2)}}
\newcommand{\img}{\ensuremath{\mathrm{img}}}
\newcommand{\bddastar}{BDDA\ensuremath{^{*}}}
\newcommand{\astar}{A\ensuremath{^{*}}}
\newcommand{\ghseta}{GHSETA\ensuremath{^{*}}}
\newcommand{\symba}{SymBA\ensuremath{^{*}}}

\theoremstyle{plain}
\newtheorem{theorem}{Theorem}
\newtheorem{lemma}{Lemma}
\newtheorem{proposition}{Proposition}
\newtheorem{corollary}{Corollary}
\theoremstyle{definition}
\newtheorem{definition}{Definition}
\newtheorem{example}{Example}
\theoremstyle{remark}
\newtheorem{remark}{Remark}

\title{Heuristics with Bounded Cofactor Width for Symbolic Search}
\author{Anonymous for review}
\date{}

\begin{document}
\maketitle

% ============================== abstract =================================
\begin{abstract}
Symbolic search with binary decision diagrams (BDDs) is a leading
approach to cost-optimal classical planning, but its dominant
instantiation remains blind: partitioning the search frontier by
heuristic value can increase BDD sizes so much that even the perfect
heuristic deteriorates performance exponentially. We give a
representation-side sufficient condition for safety. A heuristic has \emph{cofactor width}
$W$ under the search's variable order if its value is computable by
scanning the state variables once in that order while distinguishing at
most $W$ cases. We prove that partitioning \emph{any} state set by such
a heuristic has total BDD size at most $W^2n$ times the augmented
input size $|B|+2$, that symbolic \astar\ with a
consistent width-$W$ heuristic incurs at most this polynomial overhead
over blind forward search, that the known exponential deterioration
requires exponential width, and that an explicit family realizes a
quadratic overhead ratio in the width. Integer potential heuristics, pattern
databases and linear merge-and-shrink abstractions instantiate the class
with widths controlled by their standard parameters. A refined cost
model distinguishes an additional cost invisible to the node bound:
product-at-evaluation $f$-ordering can multiply high-level image calls
by the number of heuristic values, while \emph{pruning-only} search
keeps blind layer structure and enjoys a bound a factor $W$ tighter.
Archived SymK runs show small successor-batch DagSize ratios and broad
coverage gains from merge-and-shrink ordering, but an audit also finds
that their legacy node convention, fragmented blind frontiers and
signed-potential implementation do not directly validate the formal
effort or call model. We therefore treat bounded width as a safety
constraint rather than a complete algorithm selector and identify
terminal-capped abstractions, arbitrary goal-directed PDBs,
delta-integrated search and contour batching as the alternatives that
require controlled comparison.
\end{abstract}

% ============================ introduction ===============================
\section{Introduction}

A common approach to solving classical planning tasks optimally is
\astar\ search \citep{hart-et-al-ieeessc1968} with an admissible heuristic
\citep{pearl-1984}. Symbolic search
\citep{mcmillan-1993,edelkamp-helmert-ecp1999} replaces the expansion of
individual states with the expansion of state sets, compactly represented
as binary decision diagrams \citep[BDDs;][]{bryant-ieeecomp1986}. Symbolic
bidirectional blind search is among the strongest algorithms for
cost-optimal planning
\citep{torralba-et-al-aij2017,speck-et-al-jair2025}, and a key component
of the winning entry of the optimal track of the International Planning
Competition (IPC) 2023.

Combining the two paradigms has a long history. \bddastar\
\citep{edelkamp-reffel-ki1998} represents a heuristic $h$ as one BDD per
heuristic value and splits each search layer into buckets of states that
share a $g$- and an $h$-value. SetA\ensuremath{^{*}} and \ghseta\
\citep{jensen-et-al-aij2008} instead fold the heuristic-value change into
the transition relations, and ADDA\ensuremath{^{*}}
\citep{hansen-et-al-sara2002} stores $h$ as an algebraic decision diagram
(ADD). Empirically, all of these variants beat blind symbolic search in
some domains and lose in others
\citep{torralba-phd2015,torralba-et-al-aij2017}.
\citet{speck-et-al-icaps2020} explained why: in symbolic search the
effort of an expansion is governed by the number of BDD nodes rather than
the number of represented states, and they prove that even the perfect
heuristic \hstar\ can increase the cumulative size of the expanded BDDs
exponentially. Consequently, heuristic symbolic search lost traction, the
one notable exception being operator-potential heuristics
\citep{fiser-et-al-aaai2022,fiser-et-al-aij2024}, which integrate
LP-based potential heuristics \citep{pommerening-et-al-aaai2015} into
\ghseta\ and empirically ``hardly ever suffer'' from increased BDD sizes.
Why this particular family behaves so well has remained without formal
explanation, and no condition is known under which a heuristic
\emph{provably} cannot hurt symbolic search. Recently,
\citet{speck-helmert-ecai2025} gave the first worst-case performance
guarantees for symbolic search relative to explicit search, but their
analysis concerns the representation of transition relations and goal
formulas and explicitly leaves heuristics aside.

In this paper we close this gap. We introduce \emph{width-bounded
heuristics}: a heuristic has cofactor width $W$ under a variable order
if its value can be determined by reading the state variables once in
that order while maintaining at most $W$ distinct cases at any point,
i.e., if every cut through its quasi-reduced ADD in the search's variable
order contains at most $W$ distinct residual functions. (A reduced ADD may
skip levels, so its displayed node count at one level is not equivalent.) Our
contributions are as follows.
First, we prove a \emph{fragmentation theorem}: partitioning an arbitrary
state-set BDD by the values of a width-$W$ heuristic yields buckets whose
total size is bounded by a polynomial in $W$ times the size of the
original BDD---for every state set, every layer and every task. The bound
covers both known failure modes of heuristic symbolic search, the
fragmentation of the frontier and the fact that a subset of a state set
can have an exponentially larger BDD. Second, we prove a
\emph{search-effort theorem}: \bddastar\ with a consistent width-$W$
heuristic expands BDDs whose cumulative size exceeds that of blind
symbolic forward search by at most a factor $O(W^{2}n)$, where $n$ is
the number of binary variables. Third, we delimit the bound from both sides: the pathological
heuristic in the construction of \citet{speck-et-al-icaps2020} has width
$2^{\Omega(n)}$, so exponential width is necessary for their exponential
deterioration, and an explicit family of tasks and width-$W$ heuristics
realizes an $\Omega(W^2)$ overhead. Because that family has
$n=\Omega(W)$, it leaves open whether the joint $W^2n$ upper bound can
be improved. Fourth, we identify natural width-bounded families:
integer potential heuristics, pattern databases with a bounded abstract
state space, and merge-and-shrink heuristics with a linear merge
strategy---whose shrink size limit \emph{is} a width bound. This yields a
representation-effort guarantee for integer potential level sets;
delta-partitioned operator-potential search requires a separate analysis. Fifth, we refine the cost model of symbolic
search to count high-level image operations, which the node-based effort
measure cannot see, and analyze a \emph{pruning-only} search that
keeps the layer structure of blind search, uses the heuristic solely for
provably sound discards, and improves the effort bound by a factor $W$.
Finally, we implement all of the above in the SymK planner
\citep{speck-et-al-jair2025} and evaluate it on the IPC benchmarks: the
archived coverage data show that merge-and-shrink $f$-ordering improves
on blind forward search in 21 of 46 domains, while a candid audit
delimits what the old node logs and product-at-evaluation potential runs
can establish. This motivates corrected instrumentation and a
controlled comparison of terminal caps, goal-directed patterns,
delta-integrated search and image-batched contours rather than a claim
that pruning-only use is best.

% ============================ background =================================
\section{Background}

\subsection{Classical Planning}

We consider planning tasks in the SAS$^{+}$ formalism
\citep{backstrom-nebel-compint1995}.

\begin{definition}[planning task]
A \emph{planning task} is a tuple
$\task = \langle \vars, \ops, \sinit, \sgoal \rangle$, where \vars\ is a
finite set of \emph{state variables} $v$, each with a finite domain
$D_v$, \ops\ is a finite set of \emph{operators}
$o = \langle \pre_o, \eff_o \rangle$ given by partial variable
assignments and each associated with a cost
$\cost(o) \in \mathbb{Z}_{\geq 1}$, \sinit\ is a state (a total
assignment) and \sgoal\ is a partial assignment.
\end{definition}

Operator $o$ is applicable in state $s$ iff $s \models \pre_o$, and
applying it yields the state $s[o]$ that agrees with $\eff_o$ where
defined and with $s$ elsewhere. A \emph{plan} for $s$ is a sequence of
operators leading from $s$ to a goal state; its cost is the sum of the
operator costs. The perfect heuristic $\hstar(s)$ is the cost of a
cheapest plan for $s$, or $\infty$ if none exists, and
$\gstar(s)$ is the cost of a cheapest path from \sinit\ to $s$. We write
$\copt = \hstar(\sinit)$ and assume $0 < \copt < \infty$. A
\emph{heuristic} is a function
$h \colon \states \to \mathbb{Z}_{\geq 0} \cup \{\infty\}$. It is
\emph{admissible} if $h \leq \hstar$ and \emph{consistent} if
$h(s) \leq h(s[o]) + \cost(o)$ for all $s$ and applicable $o$, and
$h(s) = 0$ for all goal states. The blind heuristic is
$\hblind \equiv 0$. Restricting costs to positive integers excludes
zero-cost operators; this keeps the layer structure of uniform-cost
search simple and follows \citet{speck-et-al-icaps2020}, who assume unit
costs.

\subsection{Binary Decision Diagrams}

We fix a binary encoding of states: each variable $v$ is encoded by
$\lceil \log_2 |D_v| \rceil$ Boolean variables. Each heuristic is given
a fixed total extension to Boolean assignments that encode no state;
the extension need not be the same across heuristic families, but its
cofactor width is part of the representation being analyzed
(Remark~\ref{rem-encoding}). Let
$x_1, \dots, x_n$ denote the Boolean variables in a fixed order $\pi$; we
call $\pi$ \emph{variable-contiguous} if the bits of each SAS$^{+}$
variable are consecutive. A set of states $S$ is represented by its
characteristic function $\chi_S \colon \{0,1\}^n \to \{0,1\}$ as a
reduced ordered BDD $\bdd{S}$ with respect to $\pi$
\citep{bryant-ieeecomp1986}. We write $\bddsize{S}$ for the number of inner
nodes of $\bdd{S}$. For a Boolean or integer-valued function $f$ over
$x_1, \dots, x_n$ and a prefix assignment $\rho \in \{0,1\}^i$ to
$x_1, \dots, x_i$, the \emph{cofactor} $f|_\rho$ is the function of
$x_{i+1}, \dots, x_n$ obtained by fixing the prefix. We write
$\sub{i}{f} = \{\, f|_\rho \mid \rho \in \{0,1\}^i \,\}$ for the set of
distinct cofactors of $f$ after $i$ variables, with
$\sub{0}{f} = \{f\}$. Reduced ordered BDDs are canonical: the nodes of
$\bdd{f}$ labeled $x_{i+1}$ correspond exactly to the cofactors in
$\sub{i}{f}$ that essentially depend on $x_{i+1}$
\citep{wegener-2000}. The following quantity, which counts all cofactors
including those represented implicitly by long edges and constants, is
central to our analysis.

\begin{definition}[cofactor count]
\label{def-cofcount}
For a function $f$ over $x_1, \dots, x_n$, the \emph{cofactor count} is
$\cofcount(f) = \sum_{i=0}^{n-1} |\sub{i}{f}|$.
\end{definition}

$\cofcount(f)$ equals the number of inner nodes of the quasi-reduced
(complete) OBDD of $f$, in which every root-to-sink path tests every
variable. Both size measures are polynomially related:

\begin{lemma}
\label{lem-qsize}
For every Boolean $f$ over $x_1, \dots, x_n$,
$\bddsize{f} \leq \cofcount(f) \leq n \cdot (\bddsize{f} + 2)$.
\end{lemma}

\begin{proof}
For the first inequality, consider a node $N$ of $\bdd{f}$ labeled
$x_{i+1}$. Since every node of a reduced OBDD is reachable from the root,
there is a path from the root to $N$; fixing the variables tested along
this path to the corresponding edge values and all remaining variables
among $x_1, \dots, x_i$ arbitrarily yields a prefix $\rho \in \{0,1\}^i$
with $f|_\rho$ equal to the function computed by $N$, because variables
skipped by long edges do not influence the traversal. Hence every node
labeled $x_{i+1}$ computes an element of $\sub{i}{f}$, and by canonicity
distinct nodes compute distinct functions, so the map is injective.
Summing over $i$ gives
$\bddsize{f} \leq \sum_{i=0}^{n-1} |\sub{i}{f}| = \cofcount(f)$. For the
second inequality, fix $i$ and consider any $u \in \sub{i}{f}$. Either
$u$ is a constant function, of which there are at most two, or $u$
essentially depends on some variable; let $x_j$ with $j > i$ be the first
such variable. Traversing $\bdd{f}$ from the root under any prefix
witnessing $u$ ends at an inner node labeled $x_j$ that computes $u$, and
by canonicity this node is unique. Hence
$|\sub{i}{f}| \leq \bddsize{f} + 2$ for every $i$, and summing over the
$n$ levels yields the claim.
\end{proof}

The factor $n$ in Lemma~\ref{lem-qsize} is unavoidable: for
$f = x_1$ we have $\bddsize{f} = 1$ but $\cofcount(f) = 2n - 1$, since
the constants $0$ and $1$ are cofactors at every level $i \geq 1$.

\subsection{Symbolic Heuristic Search and Expansion Size}

Symbolic uniform-cost search maintains, for increasing $g$, the BDD of
the layer $L_g = \{\, s \mid \gstar(s) = g \,\}$, obtained by an image
operation on the previous layers followed by removal of the closed
states \citep{torralba-et-al-aij2017}. \bddastar\
\citep{edelkamp-reffel-ki1998} additionally represents the heuristic $h$
by one BDD per value: for $v \in \img(h)$, the \emph{level set} is
$\level{h}{v} = \{\, s \mid h(s) = v \,\}$. Successor sets are
intersected with the level sets, and the resulting \emph{buckets} of
states sharing a $g$-value and an $h$-value are expanded in order of
increasing $f = g + v$. Following
\citet{kissmann-edelkamp-aaai2011,torralba-phd2015,speck-et-al-icaps2020},
ties are broken in favor of smaller $g$-values, all states with the same
$(g, v)$ are represented by a single BDD, generated states already in the
closed list are removed before bucketing, states with $h$-value $\infty$
are discarded, and the search terminates as soon as a selected bucket
contains a goal state.\footnote{Fixing the goal test to selected rather
than generated buckets matches the reading of
\citet{speck-et-al-icaps2020} that all buckets counted by
Definition~\ref{def-effort} \emph{must} be expanded before an optimal
plan is returned; they note that for their task family the choice is
immaterial.} We measure search effort exactly as
\citet{speck-et-al-icaps2020}:

\begin{definition}[expansion size, \citealp{speck-et-al-icaps2020}]
\label{def-effort}
Let $B_{\task,h}(g,v)$ denote the unique BDD expanded by \bddastar\ with
$g$-value $g$ and $h$-value $v$ on task \task\ with heuristic $h$, with
$B_{\task,h}(g,v)$ representing the empty set if no such expansion
occurs. The \emph{expansion size} of \bddastar\ is
$\effort(\task,h) = \sum_{\,g + v \leq \copt,\ g < \copt}
|B_{\task,h}(g,v)|$.
\end{definition}

For $\hblind$ all states have $h$-value $0$, and \bddastar\ coincides
with symbolic uniform-cost search, so
$\effort(\task,\hblind) = \sum_{g < \copt} \bddsize{L_g}$.
\citet{speck-et-al-icaps2020} prove that no analogue of the classical
domination result for consistent heuristics in explicit \astar\ holds:
there is a family of tasks $\task_n$ on which
$\effort(\task_n, \hstar)$ is exponentially larger than
$\effort(\task_n, \hblind)$. The reason is that a strict subset
$S' \subsetneq S$ can have an exponentially larger BDD than $S$
\citep{torralba-phd2015,speck-et-al-icaps2020}, so both splitting a
layer into buckets and pruning states from it can increase
representation size. Our goal is a property of $h$ that rules out both
effects for every state set at once.

% ===================== width-bounded heuristics ==========================
\section{Bounded Cofactor Width}

Our key definition measures how much information about the prefix of a
state a heuristic must retain in order to determine its value from the
remaining suffix. We treat a heuristic as a function
$h \colon \{0,1\}^n \to \mathbb{Z}_{\geq 0} \cup \{\infty\}$ over the
Boolean encoding.

\begin{definition}[cofactor width]
\label{def-width}
The \emph{cofactor width} of a heuristic $h$ under variable order $\pi$
is $\width_\pi(h) = \max_{0 \leq i \leq n} |\sub{i}{h}|$. We call $h$
\emph{$(\pi,W)$-width-bounded} if $\width_\pi(h) \leq W$, and write
\emph{width} for cofactor width throughout.
\end{definition}

Equivalently, $h$ is $(\pi,W)$-width-bounded iff it is computed by a
deterministic layered automaton that reads $x_1, \dots, x_n$ in the order
$\pi$ using at most $W$ states per level: the automaton states at level
$i$ are exactly the cofactors in $\sub{i}{h}$, the transition on reading
a value $b$ for $x_{i+1}$ maps the cofactor $u$ to $u|_{x_{i+1} = b}$,
and the output of a final-level state, which is a constant function, is
its value. In OBDD terminology, $\width_\pi(h)$ is the maximal number of
nodes per level of the quasi-reduced multi-terminal BDD of $h$ under
$\pi$---the classical width of a Boolean or pseudo-Boolean function
\citep{wegener-2000}, here applied to heuristics and relativized to the
order used by the search.

A word on terminology. \emph{Cofactor width} is a property of the
heuristic's \emph{representation} under a variable order and is
unrelated to the novelty width of \citet{lipovetzky-geffner-ecai2012},
which measures the size of value tuples needed to distinguish states
during search, and only indirectly related to structural graph
parameters such as pathwidth, whose boundedness is a classical
sufficient condition for small OBDDs \citep{wegener-2000}. Width-style size measures likewise underlie knowledge
compilation \citep{darwiche-marquis-jair2002,darwiche-ijcai2011}, where
bounded-width classes trade succinctness for tractable operations---the
same currency our results spend on search. Three
immediate consequences will do much of the work below.

\begin{proposition}[value range]
\label{prop-range}
If $\width_\pi(h) \leq W$, then $|\img(h)| \leq W$.
\end{proposition}

\begin{proof}
Every cofactor in $\sub{n}{h}$ is a constant function, one per attained
value, so $|\img(h)| = |\sub{n}{h}| \leq \width_\pi(h) \leq W$.
\end{proof}

We write
$\values = |h(\states) \setminus \{\infty\}|$ for the number of finite
values attained on genuine planning states, and
$T=|\img(h)|$ for the number of terminals of the chosen total
extension.  By Proposition~\ref{prop-range},
$\values\leq T\leq W$.  The distinction matters operationally:
$\values$ counts search buckets, whereas $T$ also counts any values
used only for invalid Boolean encodings.

\begin{proposition}[closure]
\label{prop-closure}
Let $h_1, h_2$ have widths $W_1, W_2$ under $\pi$ and let
$\oplus \in \{+, \max, \min\}$. Then
$\width_\pi(h_1 \oplus h_2) \leq W_1 \cdot W_2$.
\end{proposition}

\begin{proof}
For every prefix $\rho$ we have
$(h_1 \oplus h_2)|_\rho = h_1|_\rho \oplus h_2|_\rho$, so the map
$(h_1|_\rho, h_2|_\rho) \mapsto (h_1 \oplus h_2)|_\rho$ is a surjection
from a set of size at most $W_1 W_2$ onto $\sub{i}{h_1 \oplus h_2}$ for
every $i$.
\end{proof}

Hence maximization and admissible additive combinations of few
width-bounded heuristics remain width-bounded, at a multiplicative cost
per component. Safe terminal mappings provide a complementary way to
reduce the interface exposed to symbolic search.

Call a possibly signed function
$q\colon\states\to\mathbb Z\cup\{\infty\}$ a \emph{raw admissible and
consistent estimate} if $q(s)\leq h^*(s)$ for every state,
$q(s)\leq c+q(s')$ for every transition $s\xrightarrow{c}s'$, and
$q(s)\leq0$ on goals.

\begin{proposition}[safe terminal transforms]
\label{prop-transform}
Let $q$ be a raw admissible and consistent estimate, equipped with a
total extension
$q \colon \{0,1\}^{n} \to \mathbb{Z} \cup \{\infty\}$. Let
$\phi \colon \mathbb{Z} \cup \{\infty\}
 \to \mathbb{Z}_{\geq 0} \cup \{\infty\}$
be nondecreasing, satisfy $\phi(0)=0$, $\phi(\infty)=\infty$, and
\[
\phi(a+c) \leq \phi(a)+c
\qquad
\text{for every $a\in\mathbb{Z}$ and integer cost $c\geq 0$.}
\]
Then $h=\phi\circ q$ is admissible and consistent and
$\width_\pi(h)\leq\width_\pi(q)$.
\end{proposition}

\begin{proof}
If $q(s')<\infty$, consistency, monotonicity and the last condition give
\[
\phi(q(s))\leq\phi(c+q(s'))\leq c+\phi(q(s')),
\]
while if $q(s')=\infty$ the transformed consistency inequality is
immediate. Thus $h$ is consistent. Monotonicity, nonnegativity and
$\phi(0)=0$ imply
$\phi(a)=0$ for $a\leq0$, while the last condition with $a=0$ implies
$\phi(a)\leq a$ for $a\geq0$. Hence
for finite $q(s)$,
$\phi(q(s))\leq\max(0,q(s))\leq h^*(s)$. If $q(s)=\infty$,
raw admissibility implies $h^*(s)=\infty$, so transformed admissibility
also holds. Thus $h$ is admissible and zero on goals. Finally,
$(\phi\circ q)|_\rho=\phi\circ(q|_\rho)$ for every prefix $\rho$.
Terminal transformation can merge cofactors but cannot create new ones.
\end{proof}

Two useful instances are rectification
$\phi_+(a)=\max(0,a)$ and, for $K\in\mathbb{Z}_{\geq0}$,
\[
\phi_K(a)=
\begin{cases}
\min(\max(0,a),K), & a<\infty,\\
\infty, & a=\infty.
\end{cases}
\]
Thus rectifying a signed estimate and capping a nonnegative integer
heuristic to $\min(h,K)$ preserve admissibility and consistency without
increasing width. The cap has at most $K+1$ finite values, so it also
caps the number of \bddastar\ buckets and product-at-evaluation image
calls per layer while retaining recognized dead ends. This is a
candidate design principle; whether the information loss pays off is an
experimental question.

The cap is not an arbitrary discretization. Among all transforms
satisfying Proposition~\ref{prop-transform} whose finite outputs lie in
$\{0,\dots,K\}$, $\phi_K$ is pointwise strongest: the proof above gives
$\phi(a)\leq\max(0,a)$, while the codomain gives
$\phi(a)\leq K$, hence $\phi(a)\leq\phi_K(a)$. Thus no other safe
monotone post-transform with the same terminal ceiling retains a larger
estimate at any state.

Finally, width-boundedness is essentially equivalent to
the compactness of the heuristic's ADD in the search's variable order.
Let $\mathrm{ADD}_\pi(h)$ denote the reduced ordered ADD of $h$ under
$\pi$ \citep{bahar-et-al-iccad1993} and $|\mathrm{ADD}_\pi(h)|$ its
number of inner nodes.

\begin{proposition}[ADD characterization]
\label{prop-add}
$\width_\pi(h) \leq |\mathrm{ADD}_\pi(h)| + |\img(h)|$ and
$|\mathrm{ADD}_\pi(h)| \leq n \cdot \width_\pi(h)$.
\end{proposition}

\begin{proof}
For the first inequality, every cofactor in $\sub{i}{h}$ is either a
constant, of which there are $|\img(h)|$ many, or the function computed
by a unique inner node of $\mathrm{ADD}_\pi(h)$ at a level below $i$, by
canonicity of reduced ADDs. For the second, every inner node of
$\mathrm{ADD}_\pi(h)$ labeled $x_{i+1}$ computes a distinct cofactor in
$\sub{i}{h}$, so the nodes per level number at most $\width_\pi(h)$, and
there are $n$ levels.
\end{proof}

Width-bounded heuristics are thus exactly the heuristics with small ADDs
\emph{in the variable order used by the search}. This order dependence is
essential and is what separates our results from statements about the
existence of some good order: the search fixes one $\pi$ for state sets,
transition relations and heuristic alike
\citep{kissmann-hoffmann-jair2014}.

\begin{remark}
\label{rem-encoding}
Definitions over $\{0,1\}^n$ require fixing $h$ on Boolean assignments
that encode no state. All upper bounds below hold for whichever total
extension is used to build the level-set BDDs, so an implementation
should pick an extension of low width; for the families in
Section~\ref{sec-instances} the natural arithmetic or table-based
extension works. Lower bounds only ever evaluate $h$ on genuine states
and therefore hold for every extension.
\end{remark}

% ======================= fragmentation theorems ==========================
\section{Bounded Fragmentation}
\label{sec-fragmentation}

We now bound the representational cost of splitting an arbitrary state
set by heuristic values. Throughout, fix the order $\pi$ and write
$W = \width_\pi(h)$. Recall that the classical Apply construction visits
pairs of operand nodes \citep{bryant-ieeecomp1986}. Under our inner-node
convention, each operand has at most its inner nodes plus two terminal nodes,
so this gives the safe bound
$\bddsize{f \wedge g} \leq
 (\bddsize{f}+2)(\bddsize{g}+2)$.
Applied naively with $g = \chi_{\level{h}{v}}$, it bounds a single bucket by
$(\bddsize{S}+2)(\bddsize{\level{h}{v}}+2)$, which does not control the sum
over all buckets. The next two theorems give bounds whose sum over buckets
stays small. The proofs use one observation, which we isolate.

\begin{lemma}
\label{lem-pair}
Let $S \subseteq \{0,1\}^n$, let $U \subseteq \img(h)$ and let
$\rho, \rho' \in \{0,1\}^i$ be prefixes with
$\chi_S|_\rho = \chi_S|_{\rho'}$ and $h|_\rho = h|_{\rho'}$. Then
$\chi_{S \cap h^{-1}(U)}|_\rho = \chi_{S \cap h^{-1}(U)}|_{\rho'}$.
\end{lemma}

\begin{proof}
For every suffix $\sigma \in \{0,1\}^{n-i}$,
$\chi_{S \cap h^{-1}(U)}|_\rho(\sigma) = \chi_S|_\rho(\sigma) \cdot
[h|_\rho(\sigma) \in U]$, which by assumption equals
$\chi_S|_{\rho'}(\sigma) \cdot [h|_{\rho'}(\sigma) \in U]
= \chi_{S \cap h^{-1}(U)}|_{\rho'}(\sigma)$.
\end{proof}

In other words, the cofactor of a bucket is determined by the pair
consisting of the cofactor of the state set and the cofactor of the
heuristic. Counting realized pairs yields our first main result.

\begin{theorem}[bucket bound]
\label{thm-bucket}
Let $S \subseteq \{0,1\}^n$ and $U \subseteq \img(h)$. Then
\[
\bddsize{S \cap h^{-1}(U)}
\;\leq\; \cofcount(\chi_{S \cap h^{-1}(U)})
\;\leq\; W \cdot \cofcount(\chi_S)
\;\leq\; W n \, (\bddsize{S} + 2).
\]
\end{theorem}

\begin{proof}
The first and last inequalities are Lemma~\ref{lem-qsize}. For the
middle inequality, fix $i \in \{0, \dots, n-1\}$ and define the set of
\emph{realized pairs}
$R_i = \{\, (\chi_S|_\rho,\, h|_\rho) \mid \rho \in \{0,1\}^i \,\}$.
Consider the map
$\mu \colon \sub{i}{\chi_{S \cap h^{-1}(U)}} \to R_i$ defined as
follows: for $u \in \sub{i}{\chi_{S \cap h^{-1}(U)}}$, choose any
prefix $\rho$ with $u = \chi_{S \cap h^{-1}(U)}|_\rho$ and set
$\mu(u) = (\chi_S|_\rho, h|_\rho)$. The map is well defined on the
chosen representatives, and it is injective: if
$\mu(u) = \mu(u')$ via prefixes $\rho, \rho'$, then
$\chi_S|_\rho = \chi_S|_{\rho'}$ and $h|_\rho = h|_{\rho'}$, so
$u = u'$ by Lemma~\ref{lem-pair}. Hence
$|\sub{i}{\chi_{S \cap h^{-1}(U)}}| \leq |R_i|
\leq |\sub{i}{\chi_S}| \cdot |\sub{i}{h}|
\leq W \cdot |\sub{i}{\chi_S}|$.
Summing over $i = 0, \dots, n-1$ gives
$\cofcount(\chi_{S \cap h^{-1}(U)}) \leq W \cdot \cofcount(\chi_S)$.
\end{proof}

With $U = \{v\}$, Theorem~\ref{thm-bucket} bounds the size of every
single bucket $S \cap \level{h}{v}$. Beyond that, it bounds the size of
\emph{every} state set obtained by removing value slices from $S$: with
$U = \img(h) \setminus \{\infty\}$ it caps the non-dead-end subset of
$S$, and with $U = \{\, v \mid v \leq c \,\}$ it caps $f$-value-based
pruning. It therefore also controls the subset blowup that drives the
negative result of \citet{speck-et-al-icaps2020}: no subset of $S$
selected by heuristic values can outgrow $S$ by more than a factor
$Wn$, up to additive terms. Summing over all values and applying
Proposition~\ref{prop-range} bounds the entire partition.

\begin{theorem}[partition bound]
\label{thm-partition}
Let $S \subseteq \{0,1\}^n$. Then
\[
\sum_{v \in \img(h)} \bddsize{S \cap \level{h}{v}}
\;\leq\; W^{2} \cdot \cofcount(\chi_S)
\;\leq\; W^{2} n \, (\bddsize{S} + 2).
\]
\end{theorem}

\begin{proof}
By Theorem~\ref{thm-bucket} applied to $U = \{v\}$, each of the at most
$|\img(h)| \leq W$ (Proposition~\ref{prop-range}) buckets has BDD size
at most $W \cdot \cofcount(\chi_S)$; the second inequality is
Lemma~\ref{lem-qsize}.
\end{proof}

The proof deliberately uses the coarse worst-case factor $W$ once for a
bucket and once for the number of values. A shared unique table such as
CUDD's may reuse isomorphic subgraphs while the buckets are constructed,
but the effort measure sums the reduced size of each bucket separately;
we do not claim a sharper aggregate bound from such sharing.

\begin{example}
\label{ex-parity}
The factor $n$ in Theorems~\ref{thm-bucket} and~\ref{thm-partition}
stems from Lemma~\ref{lem-qsize} and cannot be removed when sizes are
measured in reduced BDDs. Let $S = \{\, s \mid s(x_1) = 1 \,\}$, so
$\bddsize{S} = 1$, and let $h(s) = \bigl(\sum_{i=1}^{n} s(x_i)\bigr)
\bmod 2$, which has width $2$ under any order. The bucket
$S \cap \level{h}{0}$ consists of the assignments with $x_1 = 1$ and an
odd number of ones among $x_2, \dots, x_n$; its reduced BDD consists of
one node for $x_1$ and a parity chain with one node for $x_2$ and two
nodes for each of $x_3, \dots, x_n$, i.e., $2n - 2$ nodes in total: a
factor $2n - 2$ larger than $\bdd{S}$ although $W = 2$. Measured in
cofactor counts, where $\cofcount(\chi_S) = 2n - 1$, the bound of
Theorem~\ref{thm-bucket} is met within a factor of about two.
\end{example}

% ======================= search-effort guarantees ========================
\section{Search-Effort Guarantees}
\label{sec-effort}

We lift the set-level bounds to a guarantee on the expansion size of
\bddastar. The first step identifies the expanded buckets with slices of
the blind layers. This correspondence is folklore for explicit \astar\
with consistent heuristics; we prove it in the symbolic bucket setting
because the completeness of buckets at expansion time, i.e., that a
bucket is never re-populated after its expansion, is what makes
Definition~\ref{def-effort} well defined and our accounting exact.

\begin{lemma}[bucket identity]
\label{lem-identity}
Let $h$ be consistent. \bddastar\ on \task\ with $h$ selects every
bucket key $(g,v)$ with $g + v \leq \copt$ and $g < \copt$ and
$L_g \cap \level{h}{v} \neq \emptyset$ exactly once before terminating,
and the BDD expanded with key $(g,v)$ represents exactly
$L_g \cap \level{h}{v}$. No other bucket is expanded before
termination.
\end{lemma}

\begin{proof}
Throughout, the \emph{key} of a generated state $s$ is the pair
$(c, h(s))$, where $c$ is the cost of the generating path, and the
$f$-value of a key $(g,v)$ is $g + v$. We prove five claims.

\emph{Claim 1 (monotone selection): if a key $(g_1,v_1)$ is selected at
step $t_1$ and a key $(g_2,v_2)$ at step $t_2 > t_1$, then
$g_1 + v_1 \leq g_2 + v_2$, and if $g_1 + v_1 = g_2 + v_2$ then
$g_1 \leq g_2$.}
For the $f$-part, every state inserted into the open list during the
expansion of a bucket with key $(g,v)$ is a successor $s[o]$ of a state
$s$ with $h(s) = v$ and has key $(g + \cost(o),\, h(s[o]))$, whose
$f$-value satisfies
$g + \cost(o) + h(s[o]) \geq g + h(s) = g + v$ by consistency. Hence
expanding a minimum-$f$ key only inserts keys of at least that
$f$-value, the minimum $f$-value of the open list never decreases, and
selected $f$-values are non-decreasing over time. For the $g$-part,
suppose for contradiction that there are selections at steps
$t_1 < t_2$ of keys with equal $f$-value $F$ and $g_2 < g_1$, and among
all such violations choose one minimizing $t_2$. The bucket of
$(g_2,v_2)$ must be empty at step $t_1$: otherwise the tie-breaking
rule would select it instead of $(g_1,v_1)$. Its content is also not
inserted at step $t_1$ itself, since states inserted by the expansion
of $(g_1,v_1)$ have $g$-component $g_1 + \cost(o) > g_1 > g_2$. So
every state of the bucket is inserted at some step $t'$ with
$t_1 < t' < t_2$ by expanding a key $(g',v')$ with
$g' = g_2 - \cost(o) < g_2$ for some operator $o$. Consistency gives
$g' + v' \leq g_2 + v_2 = F$ as above, and the $f$-part applied to
$t_1 < t'$ gives $g' + v' \geq F$, so $g' + v' = F$ and $g' < g_1$:
the pair $(t_1, t')$ is a violation with $t' < t_2$, contradicting
minimality.

\emph{Claim 2 (no re-population): after a key $(g,v)$ is selected at
step $t$, no state with key $(g,v)$ is inserted into the open list;
consequently every key is selected at most once, and at its selection
its bucket contains every state that was ever generated with this key
and not expanded before.}
Suppose a state with key $(g,v)$ is inserted at a step $t' > t$ by
expanding a key $(g',v')$. Then $g' + v' \leq g + v$ by consistency as
in Claim~1, and $g' + v' \geq g + v$ by the $f$-part of Claim~1, so the
$f$-values are equal; but $g' = g - \cost(o) < g$ contradicts the
$g$-part. A second selection of the same key would require an
insertion after step $t$, which is excluded. Finally, a state generated
with key $(g,v)$ before step $t$ enters the bucket unless it is removed
by the closed-list test, which removes only previously expanded states,
and states are removed from the open list only when their bucket is
selected.

\emph{Claim 3 (optimal $g$-values): every state $s$ in a selected
bucket with key $(g,v)$ satisfies $\gstar(s) = g$ and $h(s) = v$.}
We use induction on the selection step; Claims~1 and~2 hold
unconditionally and may be used freely. That $h(s) = v$ holds by
construction of the level-set partition, and $\gstar(s) \leq g$ because
$s$ is reached by a path of cost $g$. Suppose $\gstar(s) < g$ for some
$s$ in the bucket selected at the current step, and fix a cheapest path
$p_0 = \sinit, \dots, p_k = s$, all of whose prefixes are cheapest
paths to their end states. Let $j$ be the smallest index such that
$p_j$ has not been expanded before the current step; $j$ exists because
$s$ qualifies, and $j \geq 1$ because \bddastar\ selects the initial
bucket $\{\sinit\}$ first. By the induction hypothesis, every earlier
expansion of $p_{j-1}$ has key $(\gstar(p_{j-1}), h(p_{j-1}))$, so
$p_j$ was generated with key $(\gstar(p_j), h(p_j))$ before the current
step. This key has not been selected yet: a selection after the
generation of $p_j$ would have expanded $p_j$ (Claim~2), contradicting
the choice of $j$, and a selection before the generation would
contradict Claim~2 directly. Hence the open list currently contains the
key $(\gstar(p_j), h(p_j))$, and consistency along the path segment
from $p_j$ to $s$ yields
$\gstar(p_j) + h(p_j) \leq \gstar(s) + h(s) < g + v$, using
$\gstar(s) < g$ and $h(s) = v$. This contradicts minimum-$f$ selection.

\emph{Claim 4 (completeness): every state $s$ with
$\gstar(s) = g < \copt$ and $g + h(s) \leq \copt$ is generated with key
$(g, h(s))$, and this key is selected before the search terminates.}
By strong induction on $g$. \emph{Generation:} for $g = 0$, $s = \sinit$
initializes the open list with key $(0, h(\sinit))$. For $g > 0$, let
$p$ be the predecessor of $s$ on a cheapest path, with
$\gstar(p) = g - \cost(o) < g$ and, by consistency,
$\gstar(p) + h(p) \leq \gstar(s) + h(s) \leq \copt$. By the induction
hypothesis, $p$ is generated with key $(\gstar(p), h(p))$ and this key
is selected before termination; by Claim~2 the bucket then contains
$p$, because any earlier expansion of $p$ would, by Claim~3, have
occurred with the same key, which is selected only once. Expanding the
bucket generates $s$ at cost $\gstar(p) + \cost(o) = g$ with key
$(g, h(s))$. \emph{Selection:} let $K = (g, h(s))$, with
$f(K) \leq \copt$ and $g < \copt$. From the step at which $K$ first
becomes nonempty, $K$ remains in the open list until it is selected
(Claim~2). The search terminates only when a selected bucket contains a
goal state, and by Claim~3 a goal state in a selected bucket has key
$(\gstar(s_G), 0)$ with $\gstar(s_G) \geq \copt$, since consistent
heuristics assign $0$ to goal states and \copt\ is the cheapest goal
cost. While $K$ is open, no such key is selected: its $f$-value is at
least $\copt \geq f(K)$, and in case of equality its $g$-component
$\gstar(s_G) \geq \copt > g$ loses the tie-break against $K$ or another
open key with smaller $g$-component. Hence the search does not
terminate while $K$ is open. Moreover, while $K$ is open every
selection picks a key with $f$-value at most $f(K) \leq \copt$; there
are at most $(\copt + 1)(\copt + 2)/2$ keys $(g', v')$ of nonnegative
integers with $g' + v' \leq \copt$, each selected at most once
(Claim~2), and the open list is nonempty while $K$ is open. Hence after
finitely many steps $K$ itself is selected, before termination.

\emph{Claim 5 (no other expansions): every bucket selected before the
terminating selection has a key $(g,v)$ with $g + v \leq \copt$ and
$g < \copt$.}
Fix a cheapest plan $q_0 = \sinit, \dots, q_r$ with $q_r$ a goal state
and all prefixes cheapest, so that consistency yields
$\gstar(q_l) + h(q_l) \leq \copt$ for every $l$. First we exhibit a
witness key: at every step at which the search has not yet terminated,
let $l$ be the smallest index such that $q_l$ is unexpanded, which
exists because $q_r$ is unexpanded before termination. If $l = 0$, the
initial key $(0, h(\sinit))$ has not been selected, since its bucket
contains $\sinit$, and is therefore open. If $l \geq 1$, then $q_{l-1}$
was expanded, by Claim~3 with key $(\gstar(q_{l-1}), h(q_{l-1}))$, so
$q_l$ was generated with key $K_l = (\gstar(q_l), h(q_l))$; exactly as
in the proof of Claim~3, $K_l$ has not been selected and is therefore
open. In all cases the open list contains a key with $f$-value at most
\copt, namely $(0, h(\sinit))$ or $K_l$; call it the witness. Second,
minimum-$f$ selection against the witness implies that every selected
key satisfies $g + v \leq \copt$. Third, a selected key with
$g + v \leq \copt$ and $g \geq \copt$ must equal $(\copt, 0)$, and we
show that its selection is the terminating one. Consider the step at
which $(\copt, 0)$ is selected and let $l$ be as above. If $l < r$,
the witness $K_l$ is open with $f$-value at most \copt\ and
$g$-component $\gstar(q_l) < \copt$; then either its $f$-value is
smaller than \copt, contradicting minimum-$f$ selection of
$(\copt, 0)$, or it equals \copt\ and the tie-breaking rule prefers
$K_l$ or another open key of smaller $g$-component over $(\copt, 0)$,
again a contradiction. Hence $l = r$, so $q_r$ was generated with key
$(\copt, 0)$ and, being unexpanded, lies in its bucket at selection
time by Claim~2. The selected bucket thus contains the goal state
$q_r$ and the search terminates. Consequently, every selection before
the terminating one has $g + v \leq \copt$ and $g < \copt$.

\emph{Assembly.} Let $(g,v)$ with $g + v \leq \copt$ and $g < \copt$
and $L_g \cap \level{h}{v} \neq \emptyset$. Every
$s \in L_g \cap \level{h}{v}$ satisfies the premise of Claim~4, so its
key $(g,v)$ is selected exactly once before termination (Claims~2
and~4), and at this selection the bucket contains $s$ (Claim~2, with
prior expansion of $s$ excluded because, by Claim~3, it would have
occurred at the same, uniquely selected key). Conversely, by Claim~3
the selected bucket contains only states with $\gstar = g$ and
$h$-value $v$. Hence the expanded BDD represents exactly
$L_g \cap \level{h}{v}$. Keys whose intersection is empty are never
selected, since buckets enter the open list only when populated, and by
Claim~5 no key outside the stated range is expanded before the search
stops.
\end{proof}

With the bucket identity in place, the effort bound follows from the
partition bound applied layer by layer.

\begin{theorem}[effort bound]
\label{thm-effort}
Let $h$ be consistent with $\width_\pi(h) \leq W$. Then
\[
\effort(\task, h)
\;\leq\; W^{2} n \cdot \bigl(\effort(\task, \hblind) + 2\copt\bigr).
\]
\end{theorem}

\begin{proof}
By Lemma~\ref{lem-identity} and Definition~\ref{def-effort},
\[
\effort(\task,h)
= \sum_{g < \copt} \;\sum_{\substack{v \in \img(h)\\ g + v \leq \copt}}
\bddsize{L_g \cap \level{h}{v}}
\;\leq\; \sum_{\substack{g < \copt\\ L_g \neq \emptyset}}
\;\sum_{v \in \img(h)}
\bddsize{L_g \cap \level{h}{v}},
\]
where the restriction to nonempty layers loses nothing because all
buckets of an empty layer are empty and the BDD of the empty set has no
inner nodes. By Theorem~\ref{thm-partition}, the inner sum is at most
$W^{2} n \,(\bddsize{L_g} + 2)$ for every $g$. Writing $N$ for the
number of nonempty layers $L_g$ with $g < \copt$ and using
$\effort(\task,\hblind) = \sum_{g < \copt} \bddsize{L_g}$, we obtain
\begin{equation}
\effort(\task,h)
\;\leq\; W^{2} n \,\bigl(\effort(\task,\hblind) + 2N\bigr),
\label{eq-effort-nonempty}
\end{equation}
and $N \leq \copt$ yields the claim.
\end{proof}

\begin{corollary}
\label{cor-effort}
If additionally every nonempty layer $L_g$ with $g < \copt$ is a proper
subset of the set of all Boolean assignments, then
$\effort(\task, h) \leq 3\, W^{2} n \cdot \effort(\task, \hblind)$.
\end{corollary}

\begin{proof}
The BDD of a nonempty proper subset has at least one inner node, so
$N \leq \effort(\task,\hblind)$ in
Equation~\ref{eq-effort-nonempty}.
\end{proof}

The condition of Corollary~\ref{cor-effort} fails only in degenerate
tasks. The state-set part of Theorem~\ref{thm-effort} transfers to
variants that visit the same semantic buckets, including Lazy
\bddastar\ at the frontier-representation level. This is not a bound on
the primitive operations of SetA\ensuremath{^{*}} or cost-general
\ghseta: those methods use transition-relation partitions, whose number
and BDD sizes require separate accounting. Two honest limitations. First, the guarantee is
one-sided by design: it caps the overhead of using $h$, while the
potential savings, namely the layers and dead ends that $h$ prunes via
the constraints $g + v \leq \copt$ and $v \neq \infty$, remain
unbounded, which is the desired asymmetry. Second, the baseline is blind
\emph{forward} search; the strongest blind configuration in practice is
bidirectional \citep{torralba-et-al-aij2017,speck-et-al-jair2025}, and
extending the analysis to bidirectional search, where the backward
frontier interacts with forward buckets, admits only the pointwise
pruning observation in Section~\ref{sec-prune}. A complete
bidirectional guarantee remains open, as does the case of bidirectional
\bddastar\ itself.

\subsection{A Quadratic Ratio in $W$ Is Realizable}

The factor $n$ in our bounds is unavoidable when width is held constant
(Example~\ref{ex-parity}). We now give a complementary family whose
heuristic-to-blind effort ratio is quadratic in $W$.

\begin{theorem}[quadratic width-ratio family]
\label{thm-tight}
For all integers $W \geq 2$ and $k \geq 4W$ there are tasks
$\task_{k,W}$ with $n = \Theta(k)$ Boolean variables and consistent,
admissible heuristics $h$ with $\width_\pi(h) = W$ such that
$\effort(\task_{k,W}, h) = \Omega(W^{2}) \cdot
\effort(\task_{k,W}, \hblind)$.
\end{theorem}

\begin{proof}
The task has binary variables $x_0,\dots,x_k$, choice variables
$v_1,\dots,v_k$, and one goal bit $t$. Initially only $x_0$ holds. For
each $i\in[k]$, two unit-cost operators advance the chain:
$o_i^0$ has precondition $x_{i-1}\wedge\neg x_i$ and effect $x_i$,
while $o_i^1$ has the same precondition and effect
$x_i\wedge v_i$. The only operator that achieves $t$ has precondition
$x_k\wedge\neg t$, effect $t$, and cost $W$. The goal is $t$, so every
solution has cost
\[
\copt=k+W.
\]
The order $\pi$ lists the $x$-block, then $v_1,\dots,v_k$, then $t$.

Let
\[
r(s)=\Bigl(\sum_{j=1}^{k}s(v_j)\Bigr)\bmod W,
\qquad
q(r)=(-r)\bmod W\in\{0,\dots,W-1\},
\]
and define
\[
h(s)=
\begin{cases}
0, & s(t)=1,\\
q(r(s)), & s(t)=0.
\end{cases}
\]
A unit-cost choice operator either leaves the residue unchanged or maps
$r$ to $r+1\bmod W$, and
$q(r)\leq1+q(r+1\bmod W)$. The goal operator can decrease $h$ by at
most $W-1$ and costs $W$. Hence $h$ is consistent. Every plan from a
non-goal state must still apply the cost-$W$ goal operator, so
$h\leq W-1\leq h^*$; thus $h$ is admissible.

The heuristic ignores the $x$-block. While the $v$-block is read, its
cofactor is determined by the partial residue, of which there are at
most $W$; after $t$ is read the cofactors are constants. All $W$ values
are attained, so Proposition~\ref{prop-range} gives the reverse
inequality and
\[
\width_\pi(h)=W.
\]

\emph{Blind effort.}
For $i\leq k$, layer $L_i$ fixes the chain prefix, fixes
$v_j=0$ for $j>i$ and $t=0$, and leaves $v_1,\dots,v_i$ free. Therefore
$\bddsize{L_i}=\Theta(k)$. There are $k+1$ such pre-goal layers and no
intermediate layers between costs $k$ and $k+W$, so
\[
\effort(\task_{k,W},\hblind)=\Theta(k^2).
\]

\emph{Heuristic effort.}
On $L_i$, $h=q(\sum_{j\leq i}v_j\bmod W)$. For $i\geq W-1$ all $W$
buckets are nonempty, and every bucket is selected because
$i+h\leq k+W-1<\copt$. Fix a layer
$i\in[\lceil3k/4\rceil,k]$ and one bucket. For every
$W\leq j\leq i-W$, prefixes of $v_1,\dots,v_j$ with distinct partial
residues induce distinct cofactors of the modular bucket: a suffix using
at most $W-1$ of the remaining variables completes one residue to the
required bucket and not another. Each such cofactor essentially depends
on $v_{j+1}$ by the same argument. Consequently the reduced OBDD has
$W$ distinct nodes at each of the $i-2W+1=\Omega(k)$ middle levels, and
the bucket has size $\Omega(kW)$. Since $k\geq4W$, the interval
$[\lceil3k/4\rceil,k]$ contains $\Theta(k)$ layers. Summing over these
layers and their $W$ buckets gives
\[
\effort(\task_{k,W},h)
 \geq \Theta(k)\cdot W\cdot\Omega(kW)
 =\Omega(k^2W^2).
\]
Together with the blind bound, this proves the claimed
$\Omega(W^2)$ ratio.
\end{proof}

Theorem~\ref{thm-tight} shows that a quadratic ratio in $W$ is
realizable. It does not make the joint $W^2n$ upper bound tight: its
construction requires $n=\Omega(W)$ and therefore does not exclude, for
example, a bound of order $Wn$. It should also be read together with the
small measured fragmentation ratios of Section~\ref{sec-experiments}:
the construction entangles the frontier in a way the natural families
rarely do.

\subsection{Exponential Width Is Necessary for the Known Pathology}

We now revisit the family $\task_n$ of \citet{speck-et-al-icaps2020} on
which \bddastar\ with \hstar\ is exponentially worse than with \hblind.
The task has binary variables
$v_1, \dots, v_{2n}$ and $x_0, \dots, x_{2n+1}$; initially only $x_0$
is true, operators $o_i$ and $\bar o_i$ with precondition
$x_{i-1} \wedge \neg x_i$ set $x_i$ true while setting $v_i$ true or
false, respectively, and once $x_0, \dots, x_{2n}$ are all true, a goal
operator $g_i$ with precondition $x_{2n} \wedge v_i \wedge v_{n+i}$ for
some $i \in \{1, \dots, n\}$ achieves the goal
$x_{2n+1}$. The variable order is
$\pi_n = x_0, \dots, x_{2n+1}, v_1, \dots, v_{2n}$, and the exponential
blowup arises in the bucket of non-dead-end states at $g = 2n$, whose
characteristic function contains
$\bigvee_{i=1}^{n} (v_i \wedge v_{n+i})$, a function with exponential
OBDD size under $\pi_n$ \citep{kissmann-phd2012,speck-et-al-icaps2020}.

\begin{theorem}[width lower bound]
\label{thm-lb}
For the tasks $\task_n$ and orders $\pi_n$ of
\citet{speck-et-al-icaps2020},
$\width_{\pi_n}(\hstar) \geq 2^{n}$.
\end{theorem}

\begin{proof}
We exhibit $2^n$ prefixes with pairwise distinct \hstar-cofactors. For
$A \subseteq \{1, \dots, n\}$ let $\rho_A$ be the prefix that assigns
the $x$-variables the encoding of chain position $2n$, i.e.,
$x_0, \dots, x_{2n}$ true and $x_{2n+1}$ false, and assigns $v_i = 1$
iff $i \in A$ for $i \in \{1, \dots, n\}$. All $\rho_A$ fix the same
variables $x_0, \dots, x_{2n+1}, v_1, \dots, v_n$, so their cofactors
lie in $\sub{m}{\hstar}$ for the same level $m = (2n + 2) + n$. Let
$A \neq A'$ and pick without loss of generality
$i \in A \setminus A'$. Define the suffix $\sigma$ over
$v_{n+1}, \dots, v_{2n}$ by $v_{n+i} = 1$ and $v_{n+j} = 0$ for
$j \neq i$. The state $\rho_A \sigma$ satisfies
$v_i \wedge v_{n+i}$, so a goal operator is applicable and
$\hstar(\rho_A \sigma) = 1$. The state $\rho_{A'} \sigma$ satisfies
$v_j \wedge v_{n+j}$ for no $j$: for $j = i$ because $v_i = 0$ under
$\rho_{A'}$, and for $j \neq i$ because $v_{n+j} = 0$ under $\sigma$.
No other operator is applicable in either state, since every $o_i$ and
$\bar o_i$ requires $\neg x_i$ and $x_1, \dots, x_{2n}$ are all true.
Hence $\rho_{A'} \sigma$ is a dead
end and $\hstar(\rho_{A'} \sigma) = \infty \neq 1 =
\hstar(\rho_A \sigma)$, so
$\hstar|_{\rho_A} \neq \hstar|_{\rho_{A'}}$, and
$|\sub{m}{\hstar}| \geq 2^{n}$. Since the argument evaluates \hstar\
only on states, it applies to every total extension of \hstar\ to
$\{0,1\}^n$ (Remark~\ref{rem-encoding}).
\end{proof}

Theorems~\ref{thm-effort} and~\ref{thm-lb} together delimit the
phenomenon discovered by \citet{speck-et-al-icaps2020}: under
Corollary~\ref{cor-effort}'s nonempty-proper-layer condition, any
heuristic that deteriorates \bddastar\ by a superpolynomial factor over
blind search must have superpolynomial width under the search order.
Their family satisfies that condition, and its perfect heuristic indeed
has exponential width.
Width-boundedness is a sufficient condition for safety, not a necessary
one: a heuristic of large width may still happen to induce small
buckets. Note also that \hstar\ on $\task_n$ takes at most
$2n + 3$ distinct values, so a small value range does not imply small
width; the width captures how much prefix information the level sets
entangle, not how many there are.

\subsection{A Refined Cost Model: Image Calls}
\label{sec-cost}

Definition~\ref{def-effort} inherits from \citet{speck-et-al-icaps2020} a
cost model that counts the BDD nodes of the expanded sets. This captures
the representational cost of partitioning, but not its \emph{operational}
cost: every expanded bucket triggers its own image computation, and an
image operation carries overhead beyond the size of its argument---the
traversal of the transition relations and the associated cache
effects. We make this explicit.

\begin{definition}[image-aware effort]
\label{def-image-effort}
For $\alpha \geq 0$, the \emph{$\alpha$-effort} of a symbolic search run
is $\effort_\alpha = \sum_{B} (\alpha + |B|)$, where the sum ranges over
the image operations performed before termination and $B$ is the
argument of each. The parameter $\alpha$ abstracts the per-call
overhead; for $\alpha = 0$ and the expansions counted by
Definition~\ref{def-effort}, the two measures coincide.
\end{definition}

\begin{proposition}[image-call counts]
\label{prop-calls}
Let $N_\ell$ be the number of nonempty semantic blind layers $L_g$ with
$g<C^*$.  An idealized blind implementation that stores each $L_g$ in
one BDD performs $N_\ell$ counted high-level image invocations.
An implementation of \bddastar\ that invokes its cost-image routine once
for every expanded nonempty $(g,v)$ bucket performs
$\sum_g\values_g$ such invocations, where
$\values_g\leq\values\leq W$. Hence it performs at most
$W N_\ell$ calls, and at most $\values N_\ell$ when the finite
valid-state value count is known.
\end{proposition}

\begin{proof}
The idealized blind implementation computes one image per nonempty
semantic layer. Bucket-expanding \bddastar\ computes one image for every
nonempty $(g,v)$ bucket. Its active $g$-layers form a subset of the
blind layers below $C^*$ by Lemma~\ref{lem-identity}, and each has
$\values_g\leq\values$ buckets. The last inequality is
Proposition~\ref{prop-range}.
\end{proof}

Thus width does cap the call term: the $\alpha$-part of
Definition~\ref{def-image-effort} is at most $\alpha WN_\ell$. What
width cannot provide is a width-independent one-call-per-layer bound
matching blind search. The finite value count $\values$ is often a much
sharper operational parameter than a coarse ADD-derived width upper
bound.

Here a call is one invocation of the planner's high-level cost-image
routine, not each relational product that routine performs. This
accounting applies to implementations that expand an explicit $(g,v)$
bucket per invocation. It does not apply verbatim
to SetA\ensuremath{^{*}} or cost-general \ghseta. The published
operator-potential implementation partitions transition relations by
operator-cost and heuristic-delta pairs, so its relational-product count
also depends on the number and representation of delta partitions
\citep{fiser-et-al-aij2024}.
Delta-integrated search may therefore avoid expansion per absolute
heuristic value, but requires a separate analysis and remains an
untested direction here. For product-at-evaluation \bddastar, the factor
is constructible: on the family of Theorem~\ref{thm-tight},
\bddastar\ performs $\Theta(W)$ image calls on each of $\Theta(k)$ late
layers, and hence $\Theta(W)$ calls per active pre-goal layer on average. The
archived experiments below did not record relational-product counts,
and their blind frontier could remain split into several BDD pieces, so
they do not constitute a direct empirical test of this proposition.
They do show that losses of the archived product-at-evaluation
potential implementation concentrate on tasks with many attained
values, which motivates rather than proves the image-multiplier
explanation.

\subsection{Pruning-Only Search}
\label{sec-prune}

The preceding analysis motivates a
variant that keeps every semantic layer in a \emph{single} BDD, as in
the idealized blind baseline of Proposition~\ref{prop-calls}, and uses
the heuristic only to discard states that
provably cannot contribute a better plan. \emph{Pruning-only search} runs
symbolic uniform-cost search and maintains the incumbent bound $c$, the
cost of the best solution found so far ($c = \infty$ initially; symbolic
forward search discovers solutions through frontier--goal cuts before
optimality is proved). From each generated layer it discards the states
with $h(s) = \infty$ and the slice $\{\, s \mid g + h(s) \geq c \,\}$,
re-applying the current bound when the layer is expanded. Admissibility
of $h$ makes both discards sound.

\begin{lemma}[slice identity]
\label{lem-slice}
Let $h$ be consistent, and let $c_g$ denote the incumbent bound when
pruning-only search expands the layer with $g$-value $g$. Since layers
are expanded in order of increasing $g$ and the incumbent bound never
increases, the expanded layer is exactly
$L_g \cap h^{-1}(\{\, v \mid v < c_g - g \,\})$.
\end{lemma}

\begin{proof}
The expanded set is contained in the slice because the bound $c_g$ is
applied at expansion time. For the converse, let
$s \in L_g$ with $g + h(s) < c_g$ and let
$p_0 = \sinit, \dots, p_k = s$ be a cheapest path, so
$\gstar(p_j) + h(p_j) \leq \gstar(s) + h(s) < c_g$ for every $j$ by
consistency along the path. Layer $\gstar(p_j)$ was expanded no later
than layer $g$, so its bound satisfied $c_{\gstar(p_j)} \geq c_g$, and
$p_j$ survived both its generation-time and expansion-time pruning; by
induction along the path, $s$ is generated at cost $g$ and survives, so
it lies in the expanded layer.
\end{proof}

Each expanded layer is thus a single value slice of the blind layer, so
Theorem~\ref{thm-bucket} with $U = \{\, v \mid v < c_g - g \,\}$ applies
directly, and only one such set is expanded per layer:

\begin{corollary}[pruning bound]
\label{cor-prune}
Let $h$ be admissible and consistent with $\width_\pi(h) \leq W$. The
cumulative size of the layers expanded by pruning-only search is at most
$W n \cdot (\effort(\task, \hblind) + 2\copt)$, one factor $W$ less than
the \bddastar\ bound of Theorem~\ref{thm-effort}, with one image
operation per layer as in blind search. With $c = \infty$ throughout
(no solution found yet) and no dead ends, the search coincides with
blind forward search.
\end{corollary}

Pruning-only search trades the $f$-ordered expansion of \bddastar\ for
the layer structure of blind search. Without recognized dead ends it
reaches the first solution in the same semantic $g$-layer as blind search,
then prunes the remaining optimality-proof effort by the
$g + h \geq c$ slices; recognized dead ends can be discarded from the
start. In the refined cost model its
$\alpha$-effort is at most
$W n \cdot (\effort(\task, \hblind) + 2\copt) + \alpha N_\ell$ for every
$\alpha$---the same image-call term as blind search, in contrast to
Proposition~\ref{prop-calls}. It thereby preserves the one-image-per-layer schedule and supplies the
factor-$Wn$ representation bound. Removing states need not strictly
reduce BDD size---a subset can have a larger BDD---but
Theorem~\ref{thm-bucket} bounds that possible increase.

Unlike the $f$-ordered expansion of \bddastar, pruning is sound in both
directions of bidirectional blind search. Each direction uses its own
admissible estimate: the forward
frontier discards $\{\, s \mid g_{\mathrm{fw}} + h(s) \geq c \,\}$ with
$h$ estimating goal distance, the backward frontier
$\{\, s \mid g_{\mathrm{bw}} + \bar{h}(s) \geq c \,\}$ with $\bar{h}$
estimating the distance \emph{from the initial state}. For any fixed
state set and incumbent threshold, Theorem~\ref{thm-bucket} bounds the
BDD-size increase of the corresponding slice by $Wn$ up to its additive
term. We do not lift this pointwise fact to a cumulative bidirectional
effort bound: pruning changes the forward and backward closed sets and
their meeting points, which can change solution discovery, the incumbent
sequence, and an adaptive direction schedule. Merely freezing the latter
two does not identify the pruned frontiers with fixed blind reference
sets. A self-contained guarantee must account for all of this feedback
and is left for future work. A single linear merge-and-shrink abstraction
conveniently supplies both estimates---the goal distances and the init
distances of the same shrunk transition system---and, because pruning
never partitions, this variant also requires no restriction to positive
operator costs for its soundness.


% ============================ instances ==================================
\section{Width-Bounded Heuristic Families}
\label{sec-instances}

We now show that several practically important heuristic families are
width-bounded, with widths governed by natural parameters. Throughout,
$\pi$ is variable-contiguous, $d = \max_{v \in \vars} |D_v|$ and
$m = |\vars|$, so $n \leq m \lceil \log_2 d \rceil$.

\subsection{Potential Heuristics}

A potential function \citep{pommerening-et-al-aaai2015} assigns a
potential $P(v,w)\in\mathbb{R}$ to every fact and forms the raw estimate
\[
p(s)=\sum_{v\in\vars}P(v,s(v)).
\]
The linear constraints make $p$ admissible and consistent and ensure
$p\leq0$ on goals, but $p$ may be negative. The heuristic exposed to
search is therefore
\[
\hpot(s)=\max(0,p(s)).
\]
Proposition~\ref{prop-transform} shows that this rectification preserves
admissibility and consistency and cannot increase width.
\citet{fiser-et-al-aaai2022} enforce integer potentials inside the
optimization to integrate them into \ghseta; we adopt their integer
setting.

\begin{proposition}
\label{prop-pot}
Let $p$ have integer potentials with $|P(v,w)|\leq M$ for all facts,
extended to non-state assignments by the same sum over the bit
encoding's induced table, with table values on invalid bit codes also
in $[-M,M]$, e.g., $0$, and let $\pi$ be variable-contiguous. Then
\[
\width_\pi(\hpot)\leq\width_\pi(p)\leq d(2mM+1)
\qquad\text{and}\qquad
|\img(\hpot)|\leq mM+1.
\]
\end{proposition}

\begin{proof}
Consider a prefix $\rho$ that fixes the variables
$v_1,\dots,v_j$ completely and the first
$b<\lceil\log_2d\rceil$ bits of $v_{j+1}$. The cofactor $p|_\rho$ is
determined by the accumulated sum
$\sum_{j'\leq j}P(v_{j'},\rho(v_{j'}))$, an integer in $[-mM,mM]$,
and by the assignment to the $b$ read bits of $v_{j+1}$. There are at
most $2mM+1$ sums and at most $2^b\leq d$ bit prefixes, so
$\width_\pi(p)\leq d(2mM+1)$. The raw total lies in $[-mM,mM]$;
rectification maps it to $\{0,\dots,mM\}$ and cannot increase width by
Proposition~\ref{prop-transform}.
\end{proof}

\begin{corollary}
\label{cor-pot}
For rectified consistent integer potential heuristics with magnitudes at
most $M$, under the conditions of Corollary~\ref{cor-effort},
\[
\effort(\task,\hpot)
 \leq 3d^2(2mM+1)^2n\cdot\effort(\task,\hblind),
\]
a factor polynomial in the task size and $M$.
\end{corollary}

Corollary~\ref{cor-pot} gives a representation-effort guarantee when
integer potentials are compiled into level sets and used by
product-at-evaluation \bddastar. It explains why their ADD and level
sets remain compact. It is not a complete performance guarantee for
delta-partitioned \ghseta, whose relational products require the
separate accounting described in Section~\ref{sec-cost}.

\subsection{Pattern Databases}

A pattern database \citep[PDB;][]{culberson-schaeffer-compint1998,%
edelkamp-ecp2001} projects \task\ onto a pattern
$P\subseteq\vars$ and stores the abstract goal distance $h^P(s)$ of the
abstract state induced by $s$.

\begin{proposition}
\label{prop-pdb}
Let $P\subseteq\vars$ be any pattern under a variable-contiguous order
$\pi$. Extend its table to invalid encodings by ignoring non-pattern
bits and mapping an invalid pattern-variable code to a fixed sink
terminal. Let
$|\states^P|=\prod_{v\in P}|D_v|$. Then
\[
\width_\pi(h^P)\leq d\,|\states^P|.
\]
\end{proposition}

\begin{proof}
Fix a bit-prefix $\rho$. Values read from variables outside $P$ are
irrelevant to this extension. If an already read pattern code is
invalid, the cofactor is the sink constant. Otherwise it is determined
by the assignments to the fully read pattern variables and, if the
current variable belongs to $P$, by its already-read bits. The number
of valid partial abstract assignments and bit buffers is at most
$|\states^P|$; adding the sink case still fits
$d|\states^P|$ for $d\geq2$ (the degenerate $d=1$ encoding has no
invalid bit code). Prefixes producing the same case have identical
completions. Once every pattern variable has been read, the cofactor is
a constant PDB value or the fixed sink, again within the stated bound.
\end{proof}

The bound does not require the pattern variables to form a prefix of
$\pi$: intervening non-pattern variables are simply skipped by the PDB
function. The abstract-state budget used by standard pattern selection
is therefore also a cofactor-width budget for arbitrary, including
goal-directed, patterns. The bound is coarse; the exact width is the
maximal cofactor count at a level of the corresponding quasi-reduced
PDB ADD.  A reduced ADD may skip levels, so its displayed node count at
a level is not the exact cofactor width; Proposition~\ref{prop-add}
instead supplies an ADD-derived upper bound. By
Proposition~\ref{prop-closure}, maximizing or admissibly adding a few
such PDBs, e.g., under a cost partitioning, multiplies their width
bounds.

\subsection{Linear Merge-and-Shrink}

Merge-and-shrink (M\&S) abstractions
\citep{helmert-et-al-jacm2014} iteratively merge factors of a factored
transition system and shrink intermediate abstractions to at most $N$
abstract states. With a \emph{linear} merge strategy, the abstraction
mapping is representable by cascading tables, equivalently a factored
mapping \citep{helmert-et-al-icaps2015}: a chain of tables
$\tau_1, \dots, \tau_m$ where $\tau_j$ maps the pair of the previous
intermediate abstract state and the value of variable $v_j$ to the next
intermediate abstract state.

\begin{proposition}
\label{prop-ms}
Let $h^{\alpha}$ be an M\&S heuristic built with a linear merge
strategy whose variable order agrees with the variable-contiguous
$\pi$ and with intermediate abstraction sizes at most $N$. Extend each
table with a distinguished sink for invalid variable codes. Then
$\width_\pi(h^{\alpha}) \leq d \cdot (N+1)$.
\end{proposition}

\begin{proof}
For a prefix $\rho$ fixing $v_1, \dots, v_j$ and $b$ bits of
$v_{j+1}$, the cofactor $h^{\alpha}|_\rho$ is determined by the
intermediate abstract state $\tau_j(\cdots)$ reached after processing
$v_1, \dots, v_j$ through the cascading tables, of which there are at
most $N$, or the distinguished sink, together with the $b$-bit buffer,
at most $d$ cases: the
remaining computation applies the fixed tables
$\tau_{j+1}, \dots, \tau_m$ and the final distance table to the suffix.
\end{proof}

Proposition~\ref{prop-ms} recasts a known compactness result in a new
role. Running \bddastar\ with M\&S heuristics compiled into BDDs goes
back to \citet{edelkamp-et-al-ecai2012}, and \citet{torralba-phd2015}
proved that linear M\&S heuristics have
BDD representations with at most polynomial overhead, enabling their
use in symbolic search; Proposition~\ref{prop-ms} sharpens the
representational statement to a per-level width bound and, through
Theorems~\ref{thm-partition} and~\ref{thm-effort}, converts it for the
first time into a guarantee on \emph{search behavior}: the shrink size
limit $N$, imposed anyway to keep M\&S construction tractable, doubles
as a bound on the harm the heuristic can do to the symbolic search that
uses it. Our merge-and-shrink configuration is thus deliberately the
configuration of \citet{edelkamp-et-al-ecai2012} made safe: same
algorithmic skeleton, now with an alignment condition (merge order
matching the search's variable order), a width guarantee, and---in
Section~\ref{sec-experiments}---an account of when it pays
off.\footnote{Our default configuration merges along a causal-graph
order that need not coincide with the search's Gamer order. The
corrected instrumentation therefore reports the ADD-derived upper bound
$A+T$ \emph{in the search order}; it does not measure exact cofactor
width directly. Section~\ref{sec-experiments} evaluates enforcing the
alignment condition.} In contrast, no small-width representation is known for
$h^{\max}$, $h^{\mathrm{add}}$ or LM-cut, whose values depend on
optimization problems that entangle all variables, and
Theorem~\ref{thm-lb} shows that perfect goal-distance information can
force exponential width.

% ============================ synthesis ==================================
\section{Synthesizing Width-Bounded Heuristics}
\label{sec-synthesis}

The results above invert the usual heuristic-selection problem for
symbolic search: instead of choosing the most informative heuristic and
hoping that the partition stays benign, we propose to optimize
informativeness \emph{subject to a width budget}. For potential
heuristics this is directly actionable. Following
\citet{fiser-et-al-aaai2022}, integer box constraints
$|P(v,w)|\leq M$ imply
$\width_\pi(\hpot)\leq d(2mM+1)$ by
Proposition~\ref{prop-pot}. Post-hoc rounding of individual coefficients
is not generally safe because it can violate operator consistency; any
coefficient quantization must instead be represented inside the integer
optimization constraints.

A representation-safe alternative is the terminal cap of
Proposition~\ref{prop-transform}. Applied after construction,
$\min(h,K)$ preserves admissibility and consistency, cannot increase
width, retains recognized dead ends, and limits the finite value count
to $K+1$. This suggests two concrete candidates for symbolic search.
First, cap the terminal distances of an informative merge-and-shrink
abstraction, retaining its aggregation while directly reducing the
bucket and image multiplier. Second, use goal-directed pattern selection
rather than restricting the pattern to a prefix of the BDD order;
Proposition~\ref{prop-pdb} gives the same abstract-state-based guarantee
for either choice. Both trade information for a cheaper symbolic
interface, so the theory alone does not predict a coverage improvement;
the experiments reported below do not yet establish one.

A third, algorithmic direction is to integrate heuristic deltas into
transition-relation partitions as in \ghseta. The related
FSETA\ensuremath{^{*}} view \citep{jensen-et-al-aij2008} compiles those
deltas into reduced operator costs. More generally, on transitions between
finite-valued states of a nonnegative consistent heuristic, the reduced cost
$c_h(s,o)=c(o)+h(s[o])-h(s)$ is nonnegative and turns $f$-contours into
uniform-cost layers; a symbolic realization must represent the corresponding
cost-and-delta relations and handle the zero-cost strata they can create.
These approaches may avoid an image per absolute heuristic value, but their
cost depends on delta counts and relation BDDs. We predeclare
operator-potential \ghseta\ as an external controlled baseline; the generic reduced-cost
realization remains a follow-up. Proposition~\ref{prop-calls} covers neither
one verbatim.

A fourth candidate amortizes product-at-evaluation without discarding
the heuristic order. In \emph{$f$-contour batching}, when the current
logical key has minimum value $F$ and path cost $g$, the implementation
exact-unions the same-$g$ buckets with $f\leq F+\Delta$, performs one
speculative image, and memoizes the generated slices. Goal tests,
closing and the logical key sequence still follow ordinary \bddastar.
Thus $\Delta=0$ is the original per-bucket algorithm, while larger
windows trade extra speculative successors for fewer high-level image
invocations. With a consistent heuristic and positive costs, the
speculation does not change the logical expansion order or optimal
solution cost. Unlike pruning-only search, it retains heuristic
guidance throughout the run. Its payoff is empirical: exact unions can
themselves be expensive, and a window can image buckets that ordinary
\bddastar\ would never need.

% ============================ experiments ================================
\section{Experiments}
\label{sec-experiments}

We implemented width instrumentation and all three width-bounded
heuristic families---width-capped integer potential heuristics, prefix
pattern databases and linear merge-and-shrink heuristics---in the SymK
planning system \citep{speck-et-al-jair2025}, which extends Fast Downward
\citep{helmert-jair2006}, and used the Downward Lab toolkit
\citep{seipp-et-al-zenodo2017} for running experiments on Intel Xeon Gold
6130 processors at 2.1\,GHz. Our benchmark universe consists of 1697
tasks from the optimal sequential tracks of the IPCs 1998--2023, a
subset of the suite of \citet{fiser-et-al-aij2024}. A frozen
translator-attested manifest separates two restrictions that the legacy
archive had conflated. Direct SAS scans (or an equivalent pinned
unit-cost proof for thirteen resource-heavy task translations) find
290 tasks with a zero-cost operator and 1407 with positive operator
costs. Separately, Fast Downward's default normalization introduces
axioms for the 30 \textsc{Pathways} tasks, which the implemented
abstraction heuristics reject. We therefore report the fair common
subset of 1377 tasks that are both positive-cost and axiom-free; the
\textsc{Pathways} tasks are not relabeled as zero-cost. The positive-cost
restriction is inherited from the theory: the layer structure underlying
the bucket identity (Lemma~\ref{lem-identity}) requires positive costs,
and extending the guarantees across zero-cost strata is open. The axiom
restriction is instead an implementation limitation. Two exceptions to this base are flagged where
they occur: the bidirectional pruning study uses the full suite of 1697
tasks, since pruning-only search needs no positive-cost restriction, and
repeated runs of bidirectional SymK differ by one or two solved tasks
because its direction selection is timing-based; we always compare runs
from the same sweep. We limit time to 30\,minutes and memory to
8\,GiB. The archived SymK sweeps used the planner's default
\texttt{MUTEX\_EDELETION} mode. In particular, the fragmentation
configuration did not include the
\texttt{mutex\_type=MUTEX\_NOT} override proposed in an earlier protocol
note. This does not affect applicability of
Theorem~\ref{thm-partition}, which holds for arbitrary input sets, but
the measurements are not mutex-free reproductions of complete layers
from \citet{speck-et-al-icaps2020}. All code, benchmarks and experiment
data are available in a supplementary archive.\footnote{A Zenodo DOI will be added for the
camera-ready version.}

The archived potential rows quoted below were produced by an earlier
implementation that exposed the signed raw sum $p$ directly. They
therefore do not test Corollary~\ref{cor-pot}; we retain them here only
to document the prior sweep while the rectified experiments are
regenerated.

We focus on coverage (number of solved tasks), a legacy expansion-size
diagnostic, and successor-batch fragmentation. The archived
implementation logs a partition event for each
nonempty successor batch $S$ immediately before intersection with the
heuristic level sets and reports
\[
\frac{\sum_v\operatorname{DagSize}(S\cap h^{-1}(v))}
{\operatorname{DagSize}(S)}.
\]
Here CUDD's \(\operatorname{DagSize}\) includes one constant node for
each nonempty BDD, unlike the inner-node convention
\(\bddsize{\cdot}\) of Definition~\ref{def-effort}. The two ratios are
not generally ordered because the numerator contains one such offset
per bucket. A transition-relation partition can
produce several events with the same $g$-value, so these are not
complete-layer ratios. Moreover, archived effort totals are available
only for completed runs, and the blind implementation can retain a
semantic layer as several BDD pieces after a resource-limited merge.
Coverage and runtime remain comparable, but these legacy node totals
are neither the exact metric of Definition~\ref{def-effort} nor a
theorem-faithful comparison to one BDD per blind layer. Corrected
instrumentation records inner nodes, blind piece counts and actual
high-level image invocations.

We compare against forward and bidirectional blind symbolic search in
SymK, two CPDDL \texttt{pddl-symba} configurations, and the
explicit-state cost-partitioning planner Scorpion
\citep{seipp-et-al-jair2020}. The operator-potential command uses
\texttt{--symba-fw-pot-cfg I} and
\texttt{--symba-bw-pot-cfg I}; it is therefore the I/I configuration,
not A+I. The second CPDDL configuration disables potentials and serves
as a maintained blind-bidirectional stand-in; it is not an execution of
the original IPC SymBA\ensuremath{^{*}} planner. The evaluation proceeds in six steps:
how small are the archived successor-batch ratios; how do the families
compare to blind search and the state of the art; when does $f$-ordering
beat pruning-only use of the same heuristic; what survives in
bidirectional search; what does enforcing the merge-order alignment
condition cost; and does the magnitude cap behave as a usable
accuracy--overhead knob?

\paragraph{Fragmentation.}
Across the logged partition events of solved tasks, splitting a
generated successor batch by heuristic value increases its total BDD
DagSize only marginally: the median
fragmentation ratio is $1.17$ for prefix PDBs, $1.29$ for potentials and
$1.31$ for linear merge-and-shrink, and it never exceeds $2.7$, even
though the legacy ADD diagnostic $A+\values$ ranges over five orders of
magnitude across domains (medians $141$, $576$ and $1171$, maxima up to
$5{\cdot}10^{6}$).  That diagnostic omitted a distinct terminal for
recognized dead ends, so it is not in general the width upper bound of
Proposition~\ref{prop-add}.  The corrected bound computed from a total
ADD is
\[
U=A+T,
\]
where $A$ is its number of inner nodes and $T$ its actual terminal
count.  Consequently the archived ratios are useful descriptive
evidence that the observed successor batches fragmented little, but
they do not certify that every reported run is covered by its logged
number, nor are they exact evaluations of
Theorem~\ref{thm-partition}.  The theorem's role is the worst-case
dichotomy established together with Theorems~\ref{thm-tight}
and~\ref{thm-lb}; corrected prospective runs test its metric directly.

\begin{table}[t]
\centering
\begin{tabular}{lr}
\toprule
Configuration & Coverage \\
\midrule
Blind forward (SymK) & 682 \\
Potentials ($M = 8$), \bddastar & 622 \\
Prefix PDB, \bddastar & 685 \\
Linear merge-and-shrink, \bddastar & 704 \\
Potentials ($M = 8$), pruning-only & 683 \\
Prefix PDB, pruning-only & 682 \\
Linear merge-and-shrink, pruning-only & 681 \\
\midrule
Blind bidirectional (SymK) & 777 \\
CPDDL operator potentials (I/I) \citep{fiser-et-al-aij2024} & 652 \\
CPDDL blind bidirectional stand-in & 724 \\
Scorpion \citep{seipp-et-al-jair2020} & \textbf{912} \\
\bottomrule
\end{tabular}
\caption{Coverage on the 1377 supported tasks, which are both
positive-cost and axiom-free. The upper block lists
blind forward search and our three width-bounded families, each in the
\bddastar\ (partitioning) and the pruning-only variant; the lower block
lists external baselines.}
\label{tab-coverage}
\end{table}

\paragraph{Coverage.}
Table~\ref{tab-coverage} reports total coverage and
Table~\ref{tab-domains} the per-domain breakdown. Prefix PDBs ($685$)
and linear merge-and-shrink ($704$) solve modestly more tasks than blind
forward search ($682$), whereas the archived magnitude-bounded
potential run solves substantially fewer ($622$). The width theorem
does not predict coverage or runtime: it bounds one representational
source of overhead. The merge-and-shrink family, whose abstractions
capture more global structure than a single potential function,
improves coverage, and both abstraction configurations exceed the
CPDDL I/I operator-potential configuration ($652$). This is a comparison
to the exact command executed: both directions use a single
initial-state-objective potential. It is neither an A+I run nor the
strongest multi-potential configuration reported by
\citet{fiser-et-al-aij2024}.

Per domain, the merge-and-shrink improvement is broad but individually
small: it solves more tasks than blind forward search in $21$ of $46$
domains and fewer in $5$ (sign test, $p < 0.003$), while the aggregate
difference of $22$ tasks is sensitive to domain weighting (a
domain-level bootstrap gives the wide interval $[-40, 66]$), with
\textsc{miconic} ($-23$) dominating the downside. Breadth, not aggregate
mass, is the robust part of the signal. The remaining configurations that
solve more tasks all use search paradigms outside the scope of our
forward-search guarantee (Theorem~\ref{thm-effort}): blind
\emph{bidirectional} SymK ($777$) and the CPDDL blind
bidirectional stand-in ($724$) are bidirectional,
and Scorpion ($912$) is explicit-state. Closing the gap to bidirectional
search---where the backward frontier interacts with forward buckets---is
discussed below. The robust empirical result is narrower: under this
implementation and resource limit, merge-and-shrink and the prefix PDB
add modest forward coverage, while the magnitude-bounded potential run
loses coverage. The theorem does not imply that heuristic construction,
image calls or runtime are free.

\paragraph{Ordering versus pruning.}
Table~\ref{tab-coverage} also separates two ways of using a heuristic.
The pruning-only variants finish within one task of blind forward
coverage for all three families ($683$, $682$, $681$ vs.\ $682$).
Their construction is not free: on commonly solved tasks the archived
runtime ratios are $1.02$, $1.20$ and $1.44$ for potentials, prefix
PDBs and merge-and-shrink. Ordering changes coverage in opposite
directions: linear merge-and-shrink reaches $704$ versus $681$ for
pruning-only, whereas the archived potential run reaches $622$ versus
$683$. The corresponding constructed level counts are smaller for
merge-and-shrink (median $16$) than for the raw signed potentials
(median $53$, and $65$ among observed losses).

These associations are consistent with Proposition~\ref{prop-calls},
but do not prove that value count decides the choice. The archive
contains one run per task, no call log, conditional node totals only
for completed runs, and potential rows from the pre-rectification
implementation; the blind and heuristic frontiers also use different
merge behavior. We therefore use the result to motivate a controlled
terminal-cap sweep and algorithms that reduce image invocations while
retaining guidance, not to declare pruning-only search the best
symbolic use of a heuristic.

\paragraph{Bidirectional search: pruning versus construction cost.}
Since pruning needs no partitioning, it extends to blind bidirectional
search (Section~\ref{sec-prune}), the strongest blind configuration
(1050--1051 of all 1697 tasks; repeated runs differ by a task or two).
Table~\ref{tab-bd} summarizes four successively more careful designs.
The sign never flips: eagerly built merge-and-shrink level sets cost far
more than they prune; a 60-second construction budget with fallback to
blind search, deferring construction to the first solution, and finally
seeding an incumbent from a 60-second run of \textsc{LAMA}
\citep{richter-westphal-jair2010} (exclusive bound $c$: find a cheaper
plan or prove the seed optimal; equal total budgets, coverage counts
proved-optimal outcomes only) each close most of the remaining gap, but
heuristic pruning never reliably recoups the cost of computing the
heuristic. The seeded control is the sharpest datum: the incumbent alone
lifts \emph{blind} bidirectional search to $1054$---the best
bidirectional result we observed, with no heuristic---consistent with
the fact that a bounded search need not find or reconstruct solutions.
In the forward direction the same seed completes the pruning story
positively: pruning-only potentials reach $688$ supported tasks
against $682$ for blind forward search and $683$ unseeded, the first
forward-potential configuration with a net gain. The potential for
guidance remains real but unharvested: the virtual best of blind
bidirectional search and forward merge-and-shrink solves $1081$ tasks
($30$ unique forward solves, median $947$ seconds), and neither
time-slicing nor simple pre-search selection captures it. The full study
appears in Appendix~\ref{app-bd}.

\begin{table}[t]
\centering
\begin{tabular}{lrrr}
\toprule
Configuration & Coverage & Wins & Losses \\
\midrule
Blind bidirectional & 1050--1051 & -- & -- \\
\midrule
M\&S pruning, eager & 1032 & 3 & 22 \\
\quad + 60\,s construction budget & 1041 & 3 & 13 \\
\quad + deferred to first solution & 1046 & 1 & 6 \\
\quad + seeded incumbent (\textsc{LAMA}) & 1051 & 10 & 9 \\
Seeded incumbent, no heuristic & \textbf{1054} & -- & -- \\
\bottomrule
\end{tabular}
\caption{Bidirectional pruning on the full suite (1697 tasks):
coverage and per-task wins/losses against blind bidirectional search
from the same sweep.}
\label{tab-bd}
\end{table}

\paragraph{Merge-order alignment.}
Proposition~\ref{prop-ms} requires the merge order to match the search's
variable order for the closed-form $d(N+1)$ bound; our default
merge-and-shrink configuration instead uses a causal-graph order and is
not covered by that closed-form bound. Enforcing the alignment
condition reduces the legacy ADD-derived diagnostic (median $1380$ to
$1091$, and the $90$th percentile from $50\,059$ to $9\,751$) but
costs coverage ($687$ vs.\ $704$; five tasks won,
twenty-two lost): the causal-graph merge order produces more informative
abstractions under the same shrink limit, and here accuracy outweighs
that diagnostic. Corrected $A+T$ instrumentation is needed before
interpreting the numerical change as a cofactor-width upper bound.
Width, in other words, is a constraint rather than the objective.

\paragraph{The width knob.}
Table~\ref{tab-knob} varies the potential magnitude cap $M$. At $M = 0$
the heuristic is uniformly zero and the search reduces to blind forward
search semantically, although its exact BDD-piece schedule can differ;
the legacy completed-run DagSize total is within $0.6\%$ of blind, a
useful sanity check rather than proof of implementation identity. For
$M>0$ the same conditional diagnostic has geometric-mean ratios between
$1.3$ and $1.9$, but the measurement caveats above prevent treating
these as exact expansion-size bounds. The largest observed coverage is
$623$ at $M=16$, one task above the
$622$ at $M=8$. With one run per task this does not identify a unique
optimum. We report $M=8$ as the main configuration because it has the
smaller certified magnitude bound with essentially identical observed
coverage, not because the sweep peaks there. This coefficient magnitude
bound is distinct from the post-construction terminal cap proposed in
Section~\ref{sec-synthesis}; the latter must be evaluated separately.

\begin{table}[t]
\centering
\begin{tabular}{lrrrrrrr}
\toprule
$M$ & 0 & 1 & 2 & 4 & 8 & 16 & $\infty$ \\
\midrule
Coverage & 680 & 601 & 602 & 608 & 622 & 623 & 596 \\
Legacy DagSize ratio & 1.01 & 1.92 & 1.62 & 1.45 & 1.33 & 1.33 & 1.28 \\
\bottomrule
\end{tabular}
\caption{Coverage and conditional geometric-mean legacy DagSize ratio
over blind forward search for raw potential runs with coefficient
magnitude cap $M$ ($\infty$ denotes an effectively unbounded cap).}
\label{tab-knob}
\end{table}

\begin{table}[p]
\centering
\small
\setlength{\tabcolsep}{4pt}
\begin{tabular}{lrrrrrrrrr}
\toprule
Domain & \# & b-fw & pot & PDB & M\&S & b-bd & I/I & CPDDL & Scrp \\
\midrule
airport & 50 & 23 & 22 & 23 & 23 & 24 & 27 & 27 & 37 \\
barman-opt11 & 20 & 8 & 4 & 8 & 8 & 11 & 8 & 12 & 8 \\
barman-opt14 & 14 & 3 & 3 & 3 & 3 & 6 & 3 & 6 & 3 \\
blocks & 35 & 21 & 21 & 21 & 28 & 31 & 27 & 33 & 28 \\
childsnack-opt14 & 20 & 4 & 2 & 4 & 2 & 3 & 2 & 2 & 0 \\
depot & 22 & 5 & 4 & 6 & 8 & 7 & 7 & 7 & 13 \\
driverlog & 20 & 11 & 11 & 11 & 13 & 14 & 12 & 14 & 15 \\
floortile-opt11 & 20 & 14 & 14 & 14 & 14 & 14 & 14 & 14 & 6 \\
floortile-opt14 & 20 & 20 & 17 & 20 & 17 & 20 & 20 & 20 & 5 \\
freecell & 80 & 20 & 22 & 20 & 20 & 25 & 47 & 26 & 68 \\
grid & 5 & 2 & 1 & 2 & 2 & 2 & 2 & 2 & 3 \\
gripper & 20 & 20 & 20 & 20 & 20 & 20 & 20 & 20 & 7 \\
hiking-opt14 & 20 & 16 & 15 & 16 & 16 & 18 & 13 & 18 & 19 \\
logistics00 & 28 & 16 & 20 & 17 & 21 & 20 & 0 & 0 & 28 \\
logistics98 & 35 & 5 & 5 & 5 & 6 & 6 & 5 & 5 & 12 \\
miconic & 150 & 102 & 88 & 100 & 79 & 117 & 78 & 114 & 145 \\
movie & 30 & 30 & 30 & 30 & 30 & 30 & 30 & 30 & 30 \\
mprime & 35 & 25 & 11 & 25 & 27 & 26 & 27 & 27 & 33 \\
mystery & 30 & 15 & 8 & 15 & 17 & 15 & 15 & 14 & 19 \\
nomystery-opt11 & 20 & 13 & 13 & 13 & 16 & 16 & 14 & 16 & 20 \\
organic-synthesis-opt18 & 20 & 7 & 7 & 7 & 7 & 7 & 13 & 13 & 7 \\
parking-opt11 & 20 & 0 & 0 & 0 & 4 & 1 & 1 & 2 & 8 \\
parking-opt14 & 20 & 0 & 0 & 0 & 4 & 1 & 3 & 4 & 8 \\
pipesworld-notankage & 50 & 16 & 18 & 17 & 21 & 17 & 19 & 16 & 27 \\
pipesworld-tankage & 50 & 17 & 18 & 17 & 18 & 17 & 17 & 15 & 18 \\
psr-small & 50 & 50 & 50 & 50 & 50 & 50 & 50 & 50 & 50 \\
rovers & 40 & 14 & 11 & 14 & 14 & 14 & 12 & 14 & 13 \\
satellite & 36 & 7 & 7 & 7 & 8 & 12 & 7 & 12 & 12 \\
scanalyzer-08 & 30 & 12 & 13 & 12 & 13 & 12 & 12 & 12 & 19 \\
scanalyzer-opt11 & 20 & 9 & 10 & 9 & 10 & 9 & 10 & 9 & 16 \\
snake-opt18 & 20 & 7 & 4 & 7 & 9 & 7 & 0 & 0 & 13 \\
storage & 30 & 15 & 14 & 15 & 15 & 15 & 0 & 0 & 16 \\
termes-opt18 & 20 & 12 & 12 & 14 & 12 & 18 & 13 & 18 & 7 \\
tetris-opt14 & 17 & 10 & 10 & 10 & 10 & 11 & 10 & 11 & 11 \\
tidybot-opt11 & 20 & 15 & 16 & 16 & 17 & 16 & 10 & 13 & 17 \\
tidybot-opt14 & 20 & 9 & 10 & 11 & 12 & 12 & 2 & 5 & 13 \\
tpp & 30 & 8 & 8 & 8 & 8 & 8 & 12 & 9 & 11 \\
transport-opt08 & 30 & 11 & 11 & 11 & 11 & 14 & 11 & 14 & 15 \\
transport-opt11 & 20 & 8 & 6 & 7 & 6 & 10 & 6 & 11 & 12 \\
transport-opt14 & 20 & 7 & 6 & 7 & 7 & 9 & 5 & 9 & 10 \\
trucks & 30 & 11 & 11 & 11 & 11 & 13 & 14 & 14 & 17 \\
visitall-opt11 & 20 & 12 & 13 & 12 & 13 & 12 & 12 & 12 & 17 \\
visitall-opt14 & 20 & 6 & 9 & 6 & 7 & 7 & 6 & 6 & 13 \\
woodworking-opt08 & 30 & 21 & 12 & 20 & 21 & 28 & 21 & 28 & 30 \\
woodworking-opt11 & 20 & 16 & 7 & 14 & 14 & 20 & 15 & 20 & 20 \\
zenotravel & 20 & 9 & 8 & 10 & 12 & 12 & 0 & 0 & 13 \\
\midrule
Total & 1377 & 682 & 622 & 685 & 704 & 777 & 652 & 724 & 912 \\
\bottomrule
\end{tabular}
\caption{Per-domain coverage on the 1377 supported positive-cost,
axiom-free tasks: blind forward
(b-fw), the three \bddastar\ families with bounded cofactor width
(potentials $M=8$, prefix PDB, linear merge-and-shrink), blind
bidirectional (b-bd), CPDDL operator potentials in the I/I
configuration, the CPDDL blind-bidirectional stand-in, and Scorpion
(Scrp).}
\label{tab-domains}
\end{table}

% =========================== related work ================================
\section{Related Work}

Our work connects three lines of research. First, analyses of heuristic
symbolic search: \citet{jensen-et-al-aij2008} identified the arithmetic
partitioning of frontiers as a bottleneck and mitigated it through
transition-relation branching partitions;
\citet{speck-et-al-icaps2020} proved that no heuristic, not even
\hstar, is guaranteed to help \bddastar; we give a representation-side sufficient
condition and show their pathology requires exponential width. Second,
compact representations of heuristics.
Our merge-and-shrink configuration deserves an explicit comparison,
since its algorithmic skeleton is not new:
\citet{edelkamp-et-al-ecai2012} already compiled PDB and M\&S heuristics
into vectors of BDDs and ran \bddastar, and \citet{torralba-phd2015}
proved polynomial-size BDD representations for linear M\&S and analyzed
ADD-to-BDD conversions; the factored mappings of
\citet{helmert-et-al-icaps2015}, whose cascading tables our level-set
construction walks, made the representation explicit. What none of these
works provides---and what, after the negative result of
\citet{speck-et-al-icaps2020}, made the whole approach appear unsafe---is
a bound on the effect of the representation on the \emph{search
frontier}. Our contribution to this line is therefore not the encoding
but the account of when it works: the alignment condition between merge
order and search variable order, the width bound
$\width_\pi(h^\alpha) \leq d(N+1)$ that turns the standard shrink limit
$N$ into a safety knob (Proposition~\ref{prop-ms}), and the image-call
analysis (Proposition~\ref{prop-calls}). The archived association
between this family's smaller value count and its gain from
$f$-ordering motivates the call model but, as
Section~\ref{sec-experiments} explains, does not test it directly. The operator-potential line
\citep{fiser-et-al-aaai2022,fiser-et-al-aij2024} demonstrated
empirically that potential heuristics integrate well into \ghseta.
Corollary~\ref{cor-pot} explains compact potential level sets, but does
not bound the cost-and-delta transition partitions of that
implementation. Third, performance guarantees for symbolic
search: \citet{speck-helmert-ecai2025} bound the overhead of symbolic
over explicit blind search via conjunctive transition-relation
partitioning; our results are complementary, bounding the overhead of
heuristic over blind \emph{symbolic} search via a property of the
heuristic, and combining the two guarantee regimes is an appealing
direction. Closest to our bidirectional pruning variant are the symbolic
perimeter abstraction heuristics of
\citet{torralba-et-al-ijcai2016a,torralba-et-al-aij2018}, which also
delay heuristic effort---they construct abstractions once blind search
stalls, we construct ours once the first solution makes the bound slice
applicable. The differences are the mechanism and the guarantee: SymPA
computes perimeter abstractions symbolically and uses them to order the
frontier, whereas we prune whole layers with an explicitly computed
abstraction whose init and goal distances serve both directions at once.
Theorem~\ref{thm-bucket} bounds each fixed forward or backward slice;
the cumulative bidirectional feedback remains open. A construction
budget separately bounds the setup overhead.
Related but orthogonal are variable ordering analyses
\citep{kissmann-hoffmann-jair2014} and EVMDD-based representations of
state-dependent cost structure \citep{speck-et-al-icaps2018}.

% ============================ conclusions ================================
\section{Conclusions}

We introduced heuristics of bounded cofactor width and proved that
bounded width is a sufficient condition for safety in symbolic search:
partitioning any state set by heuristic values costs at most
$W^2n$ times its augmented BDD size $|B|+2$, \bddastar\ with a
consistent width-$W$ heuristic incurs at most an $O(W^2 n)$ overhead
over blind symbolic
forward search, the known exponential deterioration requires
exponential width, and an explicit family realizes a quadratic effort
ratio in $W$ (while leaving the optimal joint dependence on $W$ and $n$
open). Potential heuristics, PDBs and linear
merge-and-shrink instantiate the class with widths controlled by the
potential magnitude, the abstract state count and the shrink limit.
Beyond the width theory, an image-call-aware cost model exposes an
operational price of product-at-evaluation $f$-ordering, and
pruning-only search carries a bound a factor $W$ tighter with the
idealized one-image-per-layer profile. The archive establishes coverage
outcomes---notably a broad merge-and-shrink gain over blind forward
search---but its CUDD-DagSize convention, completed-run-only totals,
fragmentable blind frontier, missing call log and pre-rectification
potential implementation prevent the stronger claim that it validates
the effort and call models. In bidirectional search, outside our
guarantee, the observed pruning variants do not reliably recoup
heuristic construction cost, while the incumbent alone is useful.

The audit changes the design conclusion. Bounded cofactor width is a
valuable safety constraint, but neither pruning-only use nor a small
absolute value count is by itself the best known recipe for symbolic
guidance. The immediate candidates are terminal-capped informative
abstractions, arbitrary goal-directed budget-bounded PDBs, and
$f$-contour batching, all implemented with corrected total-ADD,
inner-node, frontier-piece and image-invocation instrumentation.
Delta-integrated \ghseta\ is an essential external baseline and needs a
separate relational-product analysis. Further open directions are a
bidirectional effort guarantee that accounts for incumbent feedback and
joint optimization of abstraction, merge and BDD variable orders. We
make comparative empirical claims for the new candidates only after
their predeclared controlled evaluation.

\appendix
\section{The Bidirectional Study in Detail}
\label{app-bd}

This appendix expands the summary of Table~\ref{tab-bd}: the design
progression of the bidirectional pruning configurations and the
bound-seeding study, with the per-design diagnoses.

\paragraph{Design progression.}
Since pruning needs no partitioning, it extends to blind bidirectional
search (Section~\ref{sec-prune}), whose coverage no forward
configuration approaches. Here the abstraction itself becomes the cost:
without a construction budget, building the merge-and-shrink level sets
consumed up to the entire time limit and cost more tasks than the
pruning recovered ($1032$ vs.\ $1051$ for blind bidirectional search on
the full suite of 1697 tasks, with construction taking over $90$ seconds
on ten percent of the tasks; a tenfold larger abstraction only widens
the gap, to $990$). With a $60$-second construction budget and
fallback to blind search---making the pruning safe \emph{including} its
setup cost---the budgeted variant recovers most of the difference
($1041$ vs.\ $1051$). Deferring construction further, to the moment the
first solution makes the bound slice applicable---so that the search is
identical to blind bidirectional search up to that point---narrows the
gap again ($1046$ vs.\ $1051$; one win, six losses), but does not close
it: on the losing tasks the optimality proof consumes nearly the whole
remaining budget, and the sixty seconds invested in the abstraction
prune too little to pay for themselves. Across all three designs the
picture is consistent, and we report it as an instructive negative
result: the small legacy successor-batch DagSize ratios make
fragmentation an unlikely dominant explanation on completed runs, but
do not rule it out. At this level of performance even provably safe
pruning must first pay for computing its heuristic, and on this suite
that constant cost is not reliably recouped. The potential is
nevertheless real: the virtual best of blind bidirectional search and
the forward merge-and-shrink configuration solves $1081$ tasks---the
forward search solves $30$ tasks the bidirectional one cannot---but the
unique solves are slow (median $947$ seconds), so neither time-slicing
nor simple pre-search selection harvests the complementarity; doing so
is an open problem our data makes concrete.

\paragraph{Bound seeding.}
Since the $g + h \geq c$ slice can only prune once an incumbent exists,
we finally supplied one externally: a $60$-second run of the satisficing
planner \textsc{LAMA} \citep{richter-westphal-jair2010} yields a plan of
cost $c$, and the symbolic search
then runs with the exclusive bound $c$---it either finds a cheaper plan
or, by exhausting the bounded space, proves the incumbent optimal. All
configurations receive the same total budget, with seeding time charged
to the seeded ones, and coverage counts only proved-optimal outcomes.
The results reorder the picture in an unexpected way. Seeding helps
plain blind bidirectional search itself ($1054$ vs.\ $1050$)---consistent
with the fact that, given a bound, the search need not find or
reconstruct solutions but only raise its lower bound---so the saving
exceeds the seeding cost. Adding
merge-and-shrink pruning on top gains nothing ($1051$; it trails its own
seedless-pruning control)---the construction cost again consumes the
pruning gain. In the forward direction, by contrast, seeding completes
the pruning story positively: pruning-only potentials reach $688$
supported tasks against $682$ for blind forward search and $683$
unseeded, the first forward-potential configuration to show a net gain.
The lesson is symmetric to the bidirectional one: an incumbent is the
cheapest form of guidance---it even helps blind search---and heuristic
pruning pays off on top of it exactly when the heuristic costs
essentially nothing to build.


% ------------------------------------------------------------------------
% Bibliography from the shared aibasel/bib repo (make update-bib).
\bibliographystyle{plainnat}
\bibliography{bib/abbrv,bib/literatur,bib/crossref}
\end{document}
